\documentclass[aps,preprint,amsfonts,amssymb,showpacs,showkeys]{revtex4}%
\usepackage{amssymb}
\usepackage{amsfonts}
\usepackage{amsmath}
\usepackage{graphicx}%
\setcounter{MaxMatrixCols}{30}
%TCIDATA{OutputFilter=latex2.dll}
%TCIDATA{Version=5.50.0.2960}
%TCIDATA{CSTFile=revtex4.cst}
%TCIDATA{Created=Friday, January 07, 2011 14:59:10}
%TCIDATA{LastRevised=Monday, September 23, 2024 17:18:54}
%TCIDATA{<META NAME="GraphicsSave" CONTENT="32">}
%TCIDATA{<META NAME="SaveForMode" CONTENT="1">}
%TCIDATA{BibliographyScheme=BibTeX}
%TCIDATA{<META NAME="DocumentShell" CONTENT="Articles\SW\REVTeX 4">}
%TCIDATA{Language=American English}
%BeginMSIPreambleData
\providecommand{\U}[1]{\protect\rule{.1in}{.1in}}
%EndMSIPreambleData
\newtheorem{theorem}{Theorem}
\newtheorem{acknowledgement}[theorem]{Acknowledgement}
\newtheorem{algorithm}[theorem]{Algorithm}
\newtheorem{axiom}[theorem]{Axiom}
\newtheorem{claim}[theorem]{Claim}
\newtheorem{conclusion}[theorem]{Conclusion}
\newtheorem{condition}[theorem]{Condition}
\newtheorem{conjecture}[theorem]{Conjecture}
\newtheorem{corollary}[theorem]{Corollary}
\newtheorem{criterion}[theorem]{Criterion}
\newtheorem{definition}[theorem]{Definition}
\newtheorem{example}[theorem]{Example}
\newtheorem{exercise}[theorem]{Exercise}
\newtheorem{lemma}[theorem]{Lemma}
\newtheorem{notation}[theorem]{Notation}
\newtheorem{problem}[theorem]{Problem}
\newtheorem{proposition}[theorem]{Proposition}
\newtheorem{remark}[theorem]{Remark}
\newtheorem{solution}[theorem]{Solution}
\newtheorem{summary}[theorem]{Summary}
\newenvironment{proof}[1][Proof]{\noindent\textbf{#1.} }{\ \rule{0.5em}{0.5em}}
\begin{document}
\preprint{ }
\title[P, Q and G Conditions]{Generalized P, Q and G Conditions}
\author{Brian Weiner}
\affiliation{Department of Physics, Pennsylvania State University, DuBois PA 15801}
\keywords{N-Representability, Fock Space, Positive Extensions}
\pacs{31.15.A-,31.15.-p}

\begin{abstract}
The N-Representability Problem entails characterizing the set of second order
reduced states that are contractions of N-electron states of the Fermion Fock
algebra. This problem is formulated in the form of finding the conditions that
a positive linear functional defined on a subspace of this algebra must
satisfy in order to be extended to the whole algebra. As this algebra is a
w*-algebra one can utilize a theorem by Kadison that shows it is sufficient to
consider the values of linear functionals on projectors contained in the
subspace in order to determine whether they have positive extensions. Thus we
find the form of projectors belonging to the subspace of one and two particle
operators and subsequently show that the extension conditions needed in the
N-Representability Problem correspond to generalized P, Q and G conditions
plus the additional constraints that the functionals be dispersion free on the
number operator and their values on one particle operators determined by their
values on two particle operators.

\end{abstract}
\date{03/05/2011}
\startpage{0}
\maketitle
\tableofcontents


\section{Introduction \label{Intro}}

The $N$-Representability Problem entails finding the \emph{necessary and
sufficient conditions} that a positive trace class operator, that acts in two
electron space, must satisfy in order to be the contraction of a $N$-electron
state density operator and thus be a Second Order Reduced Density Operator
(SORDO). It was first precisely formulated in quantum chemistry by Coleman
\cite{Coleman63} as a result of Coulsons \emph{challenge }\cite{Coleman2000},
which was inspired by work initiated by Husimi \cite{Husimi}. The significance
of the problem derives from the realization that the properties of molecular
systems, at a fixed geometry, are fully described by Hamiltonians that
describe the kinetic energy of the individual electrons, interactions of the
electrons with external and internal electric and magnetic fields and
\emph{only }the simultaneous pair wise interactions of the electrons with each
other. As simultaneous three, four and higher body interactions are not
significant the $N$-electron state density operator contains redundant
information and the energy of molecular systems can be fully expressed in
terms of the SORDO of the system, which depends on fewer variables. Hence the
determination of the ground state energy can be formulated as a minimization
of a linear functional over operators that act in two electron space that are
\emph{constrained} to be SORDO's, instead of a higher dimensional minimization
over $N$-electron state density operators.

An intrinsic solution (i.e. only in terms of properties of operators acting in
two electron space) has not been found despite a large amount of effort
\cite{Mazziotti2007} along many complimentary pathways. This work has,
however, resulted in the characterization of numerous necessary conditions
(see for instance \cite{Mazziotti2007} and the references therein), the
innovative contracted Schr\"{o}dinger equation method \cite{Mazziotti} (and
the references therein), \cite{Valdemoro}, the relation to complexity theory
\cite{NPhard} and various variants of the problem i.e. electron conserving/
electron nonconserving, pure/ensemble and finite/infinite dimensional.

In a paper \cite{Weiner83} the $N$-Representability Problem was shown to be
one of finding conditions for the existence of positive extensions of a linear
functional from a \emph{subspace} of an algebra to the whole algebra. Thus
placing it firmly in the area of traditional mathematical analysis and allows
one to utilize theorems proved in this topic over the last hundred years,
where the classical Hahn-Banach theorem \cite{Encyclo} on extensions,
$\tilde{f},$ of unrestricted linear functionals, $f$, defined on subspaces,
$W,$ to the whole vector space $V$ was specialized to positive linear
functionals defined on ordered vector spaces by Krein \cite{Krein,Naimark} and
to self-adjoint subspaces of $c^{\ast}$-algebras \cite{Encyclo} that
\emph{might or might not }contain the identity by Segal \cite{Segal} and
discussed in detail in \cite{KadisonRinrose}. In order for the extension to be
a state \cite{Encyclo} i.e. a positive normalized linear functional defined on
the whole algebra its value on the identity must be one i.e.
\begin{equation}
\tilde{f}\left(  I\right)  =1
\end{equation}
so if $I\in W$ then
\begin{equation}
\tilde{f}\left(  I\right)  =f\left(  I\right)  =1.
\end{equation}
Segal's theorem examines the implications of this normalization and shows that
a positive normalized extension $\tilde{f}$ exists if (a)%

\begin{equation}
f\left(  X\right)  \geq0\,\forall X\in W_{+}, \label{Positivity}%
\end{equation}
where $W_{+}$ is the set of positive operators in $W$ i.e.
\begin{equation}
W_{+}=\left\{  X\in W\left\vert X\geq0\right.  \right\}  ,
\end{equation}
(b)%
\begin{equation}
\sup_{X\in W_{<}}f\left(  X\right)  \leq1, \label{Supcond}%
\end{equation}
where the positive part of the subspace $W$ that contains operators that are
less than or equal to the identity operator is defined as
\begin{equation}
W_{<}=\left\{  X\in W\left\vert I\geq X\geq0\right.  \right\}
\end{equation}
and (c)
\begin{equation}
\inf_{X\in W_{>}}f\left(  X\right)  \geq1, \label{Infcond}%
\end{equation}
where the positive part of the subspace $W$ that contains operators that are
greater than or equal to the identity operator is defined as
\begin{equation}
W_{>}=\left\{  X\in W\left\vert X\geq I\right.  \right\}  .
\end{equation}
It can easily be seen that \emph{if} $X\geq0$ then
\begin{equation}
\left\Vert X\right\Vert \leq1\Leftrightarrow I\geq X\geq0
\end{equation}
and thus
\begin{equation}
W_{<}=\left\{  X\in W\left\vert X\in W_{+};\left\Vert X\right\Vert
\leq1\right.  \right\}  ,
\end{equation}
and $W_{<}$ is the positive part of the unit ball
\begin{equation}
\mathfrak{B}_{1}=\left\{  X\left\vert \left\Vert X\right\Vert \leq1\right.
\right\}
\end{equation}
contained in $W$, while $W_{>}$ can be expressed as
\begin{equation}
W_{>}=\left\{  X\in W_{+}\left\vert \left\Vert X\right\Vert \geq1\right.
\right\}
\end{equation}
and is the positive part of the annulus%
\begin{equation}
\mathfrak{A}_{1}=\left\{  X\left\vert \left\Vert X\right\Vert \geq1\right.
\right\}
\end{equation}
contained in $W.$

The unit ball is a convex set and the Krein-Millmann theorem \cite{Encyclo}
shows that in a $w^{\ast}$-algebra \cite{Encyclo} it is equal to the closure
of its extreme points \cite{Kummar67}, moreover the positive part of the unit
ball in a $w^{\ast}$-algebra also forms a convex set which is also the closure
of \emph{its} extreme points. In an important theorem Kadison \cite{Kadison},
(also see the discussion in \cite{Takesaki}), has shown that the extreme
points of the positive part of the unit ball in a $w^{\ast}$-algebra are
projectors and in particular belong to the unit sphere
\begin{equation}
\left\Vert X\right\Vert =1.
\end{equation}


In this paper the subspaces we need to consider do not contain any operators
that belong to $W_{>}$ thus we only need to take into account the conditions
eqs. $\left(  \ref{Positivity}\right)  $ and $\left(  \ref{Supcond}\right)  .$
The set $W_{+}$ is also a convex set and is generated by its extreme points,
$\operatorname*{ext}\left\{  W_{+}\right\}  ,$ which is the set of projectors
in $W.$ A functional is positive on $W$ if and only if it is positive on
$\operatorname*{ext}\left\{  W_{+}\right\}  ,$ while the supremum of the
values of linear functionals defined on convex sets only occur at extreme
points of the set, thus one has only to consider the values of the functionals
on the projectors $\operatorname*{ext}\left\{  W_{+}\right\}  $ in order to
ascertain whether they satisfy both conditions eq. $\left(  \ref{Supcond}%
\right)  $ and eq. $\left(  \ref{Positivity}\right)  $. The \emph{necessary
and sufficient }conditions for the existence of a state that extends a linear
functional defined on a subspace $W$ can thus be expressed in terms of values
of the functional on the set of projectors, $\operatorname*{ext}\left\{
W_{+}\right\}  ,$ as
\begin{equation}
1\geq f\left(  X\right)  \geq0,\quad\forall X\in\operatorname*{ext}\left\{
W_{+}\right\}  . \label{NecSuff}%
\end{equation}


In the context of the $N$-Representability Problem for electrons the relevant
$w^{\ast}$-algebra is the Fermion Fock algebra and the subspace is the one
formed by one and two electron operators. In this paper the form of \emph{all
}the projectors in this subspace are characterized and the lower limits in eq.
$\left(  \ref{NecSuff}\right)  $ are shown to correspond to the well known
\cite{PQG} \emph{necessary} P, Q and G conditions. The upper limits in eq.
$\left(  \ref{NecSuff}\right)  $ correspond to the Krein, Segal and Kadison
theorems, that in conjunction with the lower limits, make the conditions
\emph{necessary and sufficient }for ensemble representability and we describe
them as \emph{generalized }P, Q and G conditions. In addition in order for a
linear functional to be $N$-Representable it must also be dispersion free on
the electron number operator and its values on the space of one particle
operators expressed in a specific manner in terms of its values on the space
of two particle operators.

In the following we discuss and examine the preceding ideas first reviewing,
in section \ref{Math}, some mathematical concepts and notation that are
utilized in the subsequent sections then in section
\ref{Expansion and Contraction} the expansion and contraction maps, which are
adjoints of each other, are considered in the context of Fock space
\cite{Kummar67,Kummar70}, which leads to the definition of Reduced Density
Operators (RDO's). In section \ref{Extensions} we utilize the fact that the
Fock algebra is a concrete realization of a $w^{\ast}$-algebra and show that
RDO's correspond to Reduced States (RS's) that are positive functionals
defined on self adjoint subspaces of a $w^{\ast}$-algebra that can be extended
to states of the algebra. This shows that the representability problems can be
formulated as finding the conditions that positive linear functionals defined
on subspaces must satisfy in order to be extended to states. Kadisons theorem
is then used to show that \emph{necessary and sufficient }extension conditions
can be expressed in terms of values of the functionals on projectors contained
in the subspace. In section \ref{Projectors} the form of \emph{all }projectors
that can be expressed as linear combinations of one and two electron operators
is found. In section \ref{PQG} we display that the conditions found from
Kadisons theorem lead to generalizations of the known \emph{necessary }P, Q
and G conditions for representability and can be expressed in terms of matrix
elements of first and second order reduced density matrices. In section
\ref{PartPair} we discuss the condition that must be satisfied to make the
functionals and reduced density matrices ensemble $N$-representable, which
leads to the list of \emph{necessary and sufficient }ensemble $N$%
-representability conditions given in section \ref{NRepCon} and we conclude
with a discussion in section \ref{Discussion}.

\section{Mathematical Essentials \label{Math}}

In the following the dimension of all spaces are considered to be finite and
can be thought of as approximations to the infinite dimensional case, which
must be used to obtain the exact solutions of quantum mechanical equations.
However notation that is appropriate for the infinite dimensional case will be
used even though it becomes redundant in the finite dimensional case.

\subsection{Fermion Fock Space \label{Fermi Fock Space}}

We consider the Fermi Fock space, $\mathcal{H}_{\mathcal{F}},$ based on the
one electron Hilbert space $\mathcal{H}^{1},$ defined by the direct sum
\begin{equation}
\mathcal{H}_{\mathcal{F}}=\underset{0\leq N\leq r}{\bigoplus}\mathcal{H}^{N}%
\end{equation}
of the $N$-electron Hilbert spaces
\begin{equation}
\mathcal{H}^{N}=\bigwedge\limits^{N}\mathcal{H}^{1},
\end{equation}
where $\wedge$ denotes antisymmetric tensor product and $\mathcal{H}%
^{0}=\left\{  \left.  \lambda\left\vert \phi\right\rangle \right\vert
\lambda\in\mathbb{C}\right\}  ,$ and $\left\vert \phi\right\rangle $ is the
vacuum state. The inner product in the Hilbert space $\mathcal{H}%
_{\mathcal{F}}$ is defined as
\begin{equation}
\left\langle \Psi\left\vert \Phi\right.  \right\rangle _{\mathcal{F}}%
=\sum_{0\leq N\leq r}\left\langle \Psi_{N}\left\vert \Phi_{N}\right.
\right\rangle _{\mathcal{H}^{N}}=\sum_{0\leq N\leq r}\left\langle
\Psi\left\vert P_{\mathcal{H}^{N}}\Phi\right.  \right\rangle _{\mathcal{H}%
^{N}},
\end{equation}
where $P_{\mathcal{H}^{N}}$ is the orthogonal projector onto $\mathcal{H}^{N}$
and $\left\langle \cdot\left\vert \cdot\right.  \right\rangle _{\mathcal{H}%
^{N}}$ is the inner product in this Hilbert space. (Notational note: If the
domain of $\left\langle \cdot\left\vert \cdot\right.  \right\rangle
_{\mathcal{H}^{N}}$ is obvious from the context the subscript "$\mathcal{H}%
^{N}"$ will not be explicitly shown.)

\subsection{Fermion Fock Algebra \label{Fermi Fock Alg}}

The fermion Fock algebra, $\mathcal{F},$ is the space of bounded operators
$\left(  \text{Appendix \ref{Operators}}\right)  $ acting in $\mathcal{H}%
_{\mathcal{F}}$ i.e.
\begin{equation}
\mathcal{F}=\mathcal{B}\left(  \mathcal{H}_{\mathcal{F}}\right)  .
\end{equation}
This space can be generated by polynomials
\begin{equation}
X=\sum_{0\leq P,Q\leq r}\sum_{\substack{1\leq j_{1}<\cdots<j_{P}\leq r\\1\leq
k_{1}<\cdots<k_{Q}\leq r}}X_{j_{1}\cdots j_{P}k_{1}\cdots k_{P}}a_{j_{1}%
}^{\dagger}\cdots a_{j_{P}}^{\dagger}a_{k_{Q}}\cdots a_{k_{1}}%
\end{equation}
of second quantized operators, $\left\{  \left.  a_{j}^{\dagger}\right\vert
1\leq j\leq r\right\}  ,$ which are defined by their action on the vacuum
vector $\left\vert \phi\right\rangle $ by
\begin{equation}
a_{j}^{\dagger}\left\vert \phi\right\rangle =\left\vert \varphi_{j}%
\right\rangle ,1\leq j\leq r,
\end{equation}
where $\left\{  \left.  \left\vert \varphi_{j}\right\rangle \right\vert 1\leq
j\leq r\right\}  $ is a complete orthonormal basis of the one-electron space
$\mathcal{H}^{1}$ and satisfy the anti-commutation relationships given by
\begin{equation}
\left[  a_{j}^{\dagger},a_{k}\right]  _{+}=a_{j}^{\dagger}a_{k}+a_{k}%
a_{j}^{\dagger}=\delta_{jk};1\leq j,k\leq r. \label{CAR}%
\end{equation}


The fermion Fock algebra can be decomposed as a vector space direct sum%
\begin{equation}
\mathcal{F=F}_{e}\mathcal{\oplus F}_{0}%
\end{equation}
of a electron conserving subalgebra $\mathcal{F}_{e}$ and a electron non
conserving subspace of operators $\mathcal{F}_{0}$. The subalgebra
$\mathcal{F}_{e}$ can be further decomposed as a direct sum of subspaces
\begin{equation}
\mathcal{F}_{e}=%
%TCIMACRO{\dbigoplus \limits_{0\leq N\leq r}}%
%BeginExpansion
{\displaystyle\bigoplus\limits_{0\leq N\leq r}}
%EndExpansion
\mathcal{B}\left(  \mathcal{H}_{\mathcal{F}}\right)  _{N}, \label{evenvect}%
\end{equation}
where
\begin{equation}
\mathcal{B}\left(  \mathcal{H}_{\mathcal{F}}\right)  _{N}=\left\{
\sum_{\substack{1\leq j_{1}<\cdots<j_{N}\leq r\\1\leq k_{1}<\cdots<k_{N}\leq
r}}X_{j_{1}\cdots j_{N}k_{1}\cdots k_{N}}a_{j_{1}}^{\dagger}\cdots a_{j_{N}%
}^{\dagger}a_{k_{N}}\cdots a_{k_{1}}\right\}
\end{equation}
or as a direct sum of subalgebras
\begin{equation}
\mathcal{F}_{e}=%
%TCIMACRO{\dbigoplus \limits_{0\leq N\leq r}}%
%BeginExpansion
{\displaystyle\bigoplus\limits_{0\leq N\leq r}}
%EndExpansion
\mathcal{B}\left(  \mathcal{H}^{N}\right)  , \label{evenalg}%
\end{equation}
where $\mathcal{B}\left(  \mathcal{H}^{N}\right)  $ is the space of bounded
operators acting in $\mathcal{H}^{N}.$ The subspace $\mathcal{F}_{0}$ can be
decomposed as the vector space direct sum
\begin{equation}
\mathcal{F}_{o}=%
%TCIMACRO{\dbigoplus \limits_{0\leq N\neq M\leq r}}%
%BeginExpansion
{\displaystyle\bigoplus\limits_{0\leq N\neq M\leq r}}
%EndExpansion
\mathcal{B}\left(  \mathcal{H}_{\mathcal{F}}\right)  _{N,M}, \label{oddsecond}%
\end{equation}
where
\begin{equation}
\mathcal{B}\left(  \mathcal{H}_{\mathcal{F}}\right)  _{N,M}=\left\{
\sum_{\substack{1\leq j_{1}<\cdots<j_{N}\leq r\\1\leq k_{1}<\cdots<k_{M}\leq
r}}X_{j_{1}\cdots j_{N}k_{1}\cdots k_{M}}a_{j_{1}}^{\dagger}\cdots a_{j_{N}%
}^{\dagger}a_{k_{M}}\cdots a_{k_{1}}\right\}
\end{equation}
or as the vector space direct sum
\begin{equation}
\mathcal{F}_{o}=%
%TCIMACRO{\dbigoplus \limits_{0\leq N\neq M\leq r}}%
%BeginExpansion
{\displaystyle\bigoplus\limits_{0\leq N\neq M\leq r}}
%EndExpansion
\mathcal{B}\left(  \mathcal{H}^{N},\mathcal{H}^{M}\right)  , \label{oddfirst}%
\end{equation}
where $\mathcal{B}\left(  \mathcal{H}^{N},\mathcal{H}^{M}\right)  $ is the
vector space of bounded linear maps from the Hilbert space $\mathcal{H}^{N}$
to the Hilbert space $\mathcal{H}^{M}.$ The decompositions eqs. $\left(
\ref{evenvect}\right)  $ and $\left(  \ref{oddsecond}\right)  $ express the
fermion Fock algebra in \emph{second quantized} form as polynomials of second
quantized operators, while the decompositions eqs. $\left(  \ref{evenalg}%
\right)  $ and $\left(  \ref{oddfirst}\right)  $ express it in \emph{first
quantized} form as expansions of the ket-bras
\begin{equation}
\left\{  \left.  \left\vert \varphi_{j_{1}}\cdots\varphi_{j_{N}}\right\rangle
\left\langle \varphi_{k_{1}}\cdots\varphi_{k_{M}}\right\vert \right\vert 1\leq
j_{1}<\cdots<j_{N}\leq r,1\leq k_{1}<\cdots<k_{M}\leq r;0\leq N,M\leq
r\right\}
\end{equation}


\subsection{Operator Inner Product \label{Inner prod}}

The Fock space trace operation is defined as
\begin{equation}
\operatorname*{Tr}\left\{  X\right\}  =\sum_{0\leq N\leq r}\sum_{1\leq
j_{1}<\cdots<j_{N}\leq r}\left\langle \Psi_{Nj_{1}\cdots j_{N}}\left\vert
X\Psi_{Nj_{1}\cdots j_{N}}\right.  \right\rangle ,
\end{equation}
where $\left\{  \left.  \Psi_{Nj_{1}\cdots j_{N}}\right\vert 1\leq
j_{1}<\cdots<j_{N}\leq r;0\leq N\leq r\right\}  $ are complete orthonormal
bases for $\mathcal{H}^{N},0\leq N\leq r.$ One can use this trace operation to
define the space of trace class operators $\mathcal{B}_{1}\left(
\mathcal{H}_{\mathcal{F}}\right)  $ that have finite trace for all bases and
the space of Hilbert Schmidt operators $\mathcal{B}_{2}\left(  \mathcal{H}%
_{\mathcal{F}}\right)  $ that form a Hilbert space with inner product
\begin{equation}
\left(  X\left\vert Y\right.  \right)  =\operatorname*{Tr}\left\{  X^{\dagger
}Y\right\}  .
\end{equation}
In finite dimensions, i.e. $r<\infty,$ there is no distinction between these
spaces of operators i.e. $\mathcal{B}\left(  \mathcal{H}_{\mathcal{F}}\right)
=\mathcal{B}_{1}\left(  \mathcal{H}_{\mathcal{F}}\right)  =\mathcal{B}%
_{2}\left(  \mathcal{H}_{\mathcal{F}}\right)  .$

\section{Expansion and Contraction Maps \label{Expansion and Contraction}}

The expansion and contraction maps relate the first and second quantization
formulation of quantum mechanics and are defined by the following:

\subsection{Contraction Map \label{Contraction}}

Employing the work of \cite{Kummar70} one can define a \emph{non singular
}generalized contraction map
\begin{equation}
L:\mathcal{B}_{1}\left(  \mathcal{H}_{\mathcal{F}}\right)  \rightarrow
\mathcal{B}_{1}\left(  \mathcal{H}_{\mathcal{F}}\right)
\end{equation}
by
\begin{equation}
L\left(  Y\right)  =\sum_{0\leq N,M\leq r}\sum_{\substack{1\leq j_{1}%
<\cdots<j_{M}\leq r\\1\leq k_{1}<\cdots<k_{N}\leq r}}\left(  \left.  a_{k_{1}%
}^{\dagger}\cdots a_{k_{N}}^{\dagger}a_{j_{M}}\cdots a_{j_{1}}\right\vert
Y\right)  \left\vert \varphi_{j_{1}}\cdots\varphi_{j_{M}}\right\rangle
\left\langle \varphi_{k_{1}}\cdots\varphi_{k_{N}}\right\vert . \label{CONT}%
\end{equation}


\subsection{Expansion Map \label{Expansion}}

The expansion map
\begin{equation}
\Gamma:\mathcal{B}\left(  \mathcal{H}_{\mathcal{F}}\right)  \rightarrow
\mathcal{B}\left(  \mathcal{H}_{\mathcal{F}}\right)
\end{equation}
is adjoint to the contraction map, i.e. $\Gamma=L^{\dagger},$ with respect to
inner product $\left(  \cdot\left\vert \cdot\right.  \right)  $ and is defined
by
\begin{equation}
\left(  \Gamma\left(  X\right)  \left\vert Y\right.  \right)  =\left(
X\left\vert L\left(  Y\right)  \right.  \right)  \quad\forall Y\in
\mathcal{B}_{1}\left(  \mathcal{H}_{\mathcal{F}}\right)  .
\end{equation}
One can show (Appendix \ref{SecQuant}) that $\Gamma$ is \emph{actually the
second quantization map}
\begin{equation}
a_{j_{1}}^{\dagger}\cdots a_{j_{N}}^{\dagger}a_{k_{P}}\cdots a_{k_{M}}%
=\Gamma\left(  \left\vert \varphi_{j_{1}}\cdots\varphi_{j_{M}}\right\rangle
\left\langle \varphi_{k_{1}}\cdots\varphi_{k_{N}}\right\vert \right)  ,
\end{equation}
which can be expressed as
\begin{equation}
\Gamma\left(  \left\vert \varphi_{j_{1}}\cdots\varphi_{j_{M}}\right\rangle
\left\langle \varphi_{k_{1}}\cdots\varphi_{k_{N}}\right\vert \right)
=\sum_{1\leq l_{1}<\cdots<l_{P}\leq r}\left\vert \varphi_{j_{1}}\cdots
\varphi_{j_{M}}\varphi_{l_{1}}\cdots\varphi_{l_{P}}\right\rangle \left\langle
\varphi_{k_{1}}\cdots\varphi_{k_{N}}\varphi_{l_{1}}\cdots\varphi_{l_{P}%
}\right\vert .
\end{equation}


The expansion and contraction maps transform naturally between the first and
second quantization direct sum decompositions of $\mathcal{F}$ i.e.
\begin{equation}
\Gamma:%
%TCIMACRO{\dbigoplus \limits_{0\leq N,M\leq r}}%
%BeginExpansion
{\displaystyle\bigoplus\limits_{0\leq N,M\leq r}}
%EndExpansion
\mathcal{B}\left(  \mathcal{H}^{N},\mathcal{H}^{M}\right)  \rightarrow%
%TCIMACRO{\dbigoplus \limits_{0\leq N,M\leq r}}%
%BeginExpansion
{\displaystyle\bigoplus\limits_{0\leq N,M\leq r}}
%EndExpansion
\mathcal{B}\left(  \mathcal{H}_{\mathcal{F}}\right)  _{N,M}%
\end{equation}
and
\begin{equation}
L:%
%TCIMACRO{\dbigoplus \limits_{0\leq N,M\leq r}}%
%BeginExpansion
{\displaystyle\bigoplus\limits_{0\leq N,M\leq r}}
%EndExpansion
\mathcal{B}\left(  \mathcal{H}_{\mathcal{F}}\right)  _{N,M}\rightarrow%
%TCIMACRO{\dbigoplus \limits_{0\leq N,M\leq r}}%
%BeginExpansion
{\displaystyle\bigoplus\limits_{0\leq N,M\leq r}}
%EndExpansion
\mathcal{B}\left(  \mathcal{H}^{N},\mathcal{H}^{M}\right)  .
\end{equation}


\subsection{Reduced Density Operators \label{Reduced Density Operators}}

The $P^{th}$ order reduced density operator associated with a state $D\in
S_{\mathcal{F}}$ can be defined by using the contraction map $L$ as
\begin{equation}
D_{P}=P_{\mathcal{H}^{P}}L\left(  D\right)  P_{\mathcal{H}^{P}}=\sum
_{\substack{1\leq j_{1}<\cdots<j_{P}\leq r\\1\leq k_{1}<\cdots<k_{P}\leq
r}}\left(  \left.  a_{k_{1}}^{\dagger}\cdots a_{k_{P}}^{\dagger}a_{j_{1}%
}\cdots a_{j_{P}}\right\vert D\right)  \left\vert \varphi_{j_{1}}\cdots
\varphi_{j_{P}}\right\rangle \left\langle \varphi_{k_{1}}\cdots\varphi_{k_{P}%
}\right\vert \label{PORDO}%
\end{equation}
and in particular the SORDO is given by
\begin{equation}
D_{2}=P_{\mathcal{H}^{2}}L\left(  D\right)  P_{\mathcal{H}^{2}}=\sum
_{\substack{1\leq j_{1}<j_{2}\leq r\\1\leq k_{1}<k_{2}\leq r}}\left(  \left.
a_{k_{1}}^{\dagger}a_{k_{2}}^{\dagger}a_{j_{1}}a_{j_{2}}\right\vert D\right)
\left\vert \varphi_{j_{1}}\varphi_{j_{2}}\right\rangle \left\langle
\varphi_{k_{1}}\varphi_{k_{2}}\right\vert .
\end{equation}
If $D$ is a state describing a system of $N$ electrons then
\begin{gather}
D=D^{N}=P_{\mathcal{H}^{N}}DP_{\mathcal{H}^{N}}\nonumber\\
P_{\mathcal{H}^{M}}DP_{\mathcal{H}^{M}}=0,M\neq N
\end{gather}
and the above construction corresponds to the normalization
\begin{equation}
\operatorname*{Tr}\left\{  D_{2}\right\}  =\binom{N}{2}.
\end{equation}


\section{Extensions of Positive Functionals \label{Extensions}}

The Fermion Fock algebra is a concrete realization of a $w^{\ast}$-algebra
\cite{Encyclo} (that contains an identity element), which is a $c^{\ast}%
$-algebra that also has the property of being a dual space of another Banach
space, in this case $\mathcal{B}\left(  \mathcal{H}_{\mathcal{F}}\right)  $ is
the dual space of $\mathcal{B}_{1}\left(  \mathcal{H}_{\mathcal{F}}\right)  $
i.e. $\mathcal{B}_{1}\left(  \mathcal{H}_{\mathcal{F}}\right)  ^{\ast}$. The
important property of $w^{\ast}$-algebras vis-\`{a}-vis $c^{\ast}$-algebras is
that they are the closure of their idempotents, which in the case of
$\mathcal{F}$ are the orthogonal projectors onto subspaces of $\mathcal{H}%
_{\mathcal{F}}.$

\subsection{States \label{States}}

A state on a $w^{\ast}$-algebra is a bounded positive linear function, $f$,
such that $f\left(  I\right)  =1$, where $I$ is the identity element of the
algebra, it is called a normal state if it can be identified with an element
of the predual of the algebra otherwise it is an abnormal state. In this paper
only normal states belonging to the predual $\mathcal{B}_{1}\left(
\mathcal{H}_{\mathcal{F}}\right)  $ of $\mathcal{B}\left(  \mathcal{H}%
_{\mathcal{F}}\right)  $ are considered. In the Fermion Fock algebra normal
states are in one to one correspondence with normalized density operators
i.e.
\begin{equation}
f\left(  X\right)  =\operatorname*{Tr}\left\{  XD_{f}\right\}  ;\,D_{f}\in
S_{\mathcal{F}}.
\end{equation}


The convex set, $S_{\mathcal{F}},$ of density operators in $\mathcal{F}$ is
given by
\begin{equation}
S_{\mathcal{F}}=\left\{  D\left\vert D\in\mathcal{B}_{1}\left(  \mathcal{H}%
_{\mathcal{F}}\right)  ;D\geq0;\operatorname{Tr}D=1\right.  \right\}  ,
\end{equation}
whose elements could describe any type of physical state i.e. electron
conserving or non conserving, ensemble or pure etc.

\subsection{Partial States \label{Partial States}}

A \emph{partial} state is a positive linear functional defined on a self
adjoint subspace, $\mathcal{W},$ of a $w^{\ast}$-algebra \emph{that contains}
the identity $I$ and is a restriction of a state. It can be shown \cite{Emch,
Krein} that the necessary and sufficient condition for a positive linear
functional defined on such a subspace $\mathcal{W}$ to be a partial state and
thus extended to a state is
\begin{equation}
\sup_{\substack{I\geq X\geq0\\X\in\mathcal{W}}}f\left(  X\right)  \leq1
\label{LessI}%
\end{equation}
and%
\begin{equation}
\inf_{\substack{X\geq I\\X\in\mathcal{W}}}f\left(  X\right)  \geq1.
\label{MoreI}%
\end{equation}


The identity element $I_{\mathcal{H}_{\mathcal{F}}}$ of $\mathcal{B}\left(
\mathcal{H}_{\mathcal{F}}\right)  $ can be expanded as
\begin{equation}
I_{\mathcal{H}_{\mathcal{F}}}=\sum_{0\leq N\leq r}P_{\mathcal{H}^{N}},
\end{equation}
where $P_{\mathcal{H}^{N}}$ are the orthogonal projectors onto $\mathcal{H}%
^{N},$ $0\leq N\leq r$ and in particular $P_{\mathcal{H}^{0}}=$ $\left\vert
\phi\right\rangle \left\langle \phi\right\vert $ is the projector onto the
vacuum state. Thus the condition eq.$\left(  \ref{MoreI}\right)  $ only
effects operators that are non zero on the vacuum state.

\subsection{Reduced States \label{Reduced States}}

A \emph{reduced} state is a postive linear functional defined on a self
adjoint subspace of a $w^{\ast}$-algebra \emph{that does not contain} the
identity and is a restriction of a state. The necessary and sufficient
condition \cite{Segal} for a linear functional defined on such a subspace
$\mathcal{W}$ to be a reduced state and thus be extended to a state is
\begin{equation}
1\geq f\left(  X\right)  \geq0,\quad\forall X\in\mathfrak{B}_{1+}\left(
\mathcal{W}\right)  ,
\end{equation}
where $\mathfrak{B}_{1+}\left(  \mathcal{W}\right)  $ is the positive part of
the unit ball of $\mathcal{B}\left(  \mathcal{H}_{\mathcal{F}}\right)  $
contained in $W$ given by%
\begin{equation}
\mathfrak{B}_{1+}\left(  \mathcal{W}\right)  =\left\{  \left.  X\right\vert
X\in\mathcal{W},I\geq X\geq0\right\}  \equiv\left\{  \left.  X\right\vert
X\in\mathcal{W}_{+},\left\Vert X\right\Vert _{\mathcal{B}\left(
\mathcal{H}_{\mathcal{F}}\right)  }\geq1\right\}  \label{BallCond}%
\end{equation}
and
\begin{equation}
\inf_{X\in\mathfrak{A}_{1+}\left(  \mathcal{W}\right)  }f\left(  X\right)
\geq1,
\end{equation}
where $\mathfrak{A}_{1+}\left(  \mathcal{W}\right)  $ is the positive part of
the unit annulus of $\mathcal{B}\left(  \mathcal{H}_{\mathcal{F}}\right)  $
contained in $W$ given by
\begin{equation}
\mathfrak{A}_{1+}\left(  \mathcal{W}\right)  =\left\{  \left.  X\right\vert
X\in\mathcal{W},X\geq I\right\}  \equiv\left\{  \left.  X\right\vert
X\in\mathcal{W}_{+},\left\Vert X\right\Vert _{\mathcal{B}\left(
\mathcal{H}_{\mathcal{F}}\right)  }\leq1\right\}  . \label{AnnCon}%
\end{equation}


An extension of an important theorem of Kadison \cite{Kadison} applied to self
adjoint subspaces of $w^{\ast}$-algebra shows that the positive extension
condition eq.$\left(  \ref{BallCond}\right)  $ can be expressed in terms of
projectors belonging $\mathcal{W},$ an observation that is crucial in deriving
necessary and sufficient conditions for representability (Section
\ref{Projectors}).

In order to obtain representability conditions for SORDO's the appropriate
subspace to consider in $\mathcal{B}\left(  \mathcal{H}_{\mathcal{F}}\right)
$ is the subspace formed by one and two electron conserving operators
\begin{equation}
\Lambda_{2}=%
%TCIMACRO{\dbigoplus \limits_{1\leq P\leq2}}%
%BeginExpansion
{\displaystyle\bigoplus\limits_{1\leq P\leq2}}
%EndExpansion
\mathcal{B}\left(  \mathcal{H}_{\mathcal{F}}\right)  _{P},
\end{equation}
as reduced states defined on $\Lambda_{2}$ correspond to SORDO's through%
\begin{equation}
D_{2}=\sum_{\substack{1\leq j_{1}<j_{2}\leq r\\1\leq k_{1}<k_{2}\leq
r}}f\left(  a_{j_{1}}^{\dagger}a_{j_{2}}^{\dagger}a_{k_{2}}a_{k_{1}}\right)
\left\vert \varphi_{j_{1}}\varphi_{j_{2}}\right\rangle \left\langle
\varphi_{k_{1}}\varphi_{k_{2}}\right\vert ,
\end{equation}
where%
\begin{equation}
f\left(  a_{j_{1}}^{\dagger}a_{j_{2}}^{\dagger}a_{k_{2}}a_{k_{1}}\right)
=\left(  \left.  a_{k_{1}}^{\dagger}a_{k_{2}}^{\dagger}a_{j_{2}}a_{j_{1}%
}\right\vert D\right)
\end{equation}
and FORDO's through
\begin{equation}
D_{1}=\sum_{\substack{1\leq j_{1}\leq r\\1\leq k_{1}\leq r}}f\left(  a_{j_{1}%
}^{\dagger}a_{k_{1}}\right)  \left\vert \varphi_{j_{1}}\right\rangle
\left\langle \varphi_{k_{1}}\right\vert ,
\end{equation}
where%

\begin{equation}
f\left(  a_{j_{1}}^{\dagger}a_{k_{1}}\right)  =\left(  \left.  a_{k_{1}%
}^{\dagger}a_{j_{1}}\right\vert D\right)  .
\end{equation}


In $\mathcal{B}\left(  \mathcal{H}_{\mathcal{F}}\right)  $ the condition
eq.$\left(  \ref{AnnCon}\right)  $ only effects operators that are non zero on
the vacuum state thus does not constrain any operators belonging to
$\Lambda_{2}$ and one can just concentrate on the conditions given by
eq.$\left(  \ref{BallCond}\right)  $.

\section{Projectors \label{Projectors}}

As we are searching for necessary and sufficient conditions for extending
positive \emph{linear} functionals the condition eq. $\left(  \ref{Supcond}%
\right)  $ needs only to be checked on extreme points of the positive portion
of the unit ball $\mathfrak{B}_{1}\left(  \Lambda_{2}\right)  ,$ which are
projection operators $\cite{Kadison}$. One should note that these conditions
will be representability conditions and that further constraints need to
considered (section \ref{NRepCon}) to obtain $N$-representability conditions.
Projectors in $\mathcal{B}\left(  \mathcal{H}_{\mathcal{F}}\right)  $ are
postive idempotent operators, $P,$ with norm equal to one i.e.
\begin{align}
P  &  \geq0\nonumber\\
P^{2}  &  =P\nonumber\\
\left\Vert P\right\Vert _{\mathcal{B}\left(  \mathcal{H}_{F}\right)  }  &
=\sup_{\left\vert \Psi\right\rangle \in\mathcal{H}_{\mathcal{F}}}\left\{
\frac{\left\langle \Psi\left\vert P\Psi\right.  \right\rangle }{\left\langle
\Psi\left\vert \Psi\right.  \right\rangle }\right\}  =1.
\end{align}
We first characterize the projectors in $\mathcal{B}\left(  \mathcal{H}%
_{\mathcal{F}}\right)  _{1},$ then $\mathcal{B}\left(  \mathcal{H}%
_{\mathcal{F}}\right)  _{2}$ and finally $\Lambda_{2}.$

\subsection{Projectors in $\mathcal{B}\left(  \mathcal{H}_{\mathcal{F}%
}\right)  _{1}$}

Every positive operator in $\mathcal{B}\left(  \mathcal{H}_{\mathcal{F}%
}\right)  _{1}$ can be expressed as (Appendix \ref{Projection} )
\begin{equation}
\lambda_{1}=\sum_{1\leq j,k\leq r}\lambda_{1jk}a_{j}^{\dagger}a_{k}%
=\sum_{1\leq l\leq r}\beta_{l}b_{l}^{\dagger}b_{l},\,\beta_{l}\geq0,1\leq
l\leq r, \label{1par}%
\end{equation}
where
\begin{equation}
b_{l}^{\dagger}=\sum_{1\leq j\leq r}u_{lj}a_{j}^{\dagger},\,1\leq l\leq r
\end{equation}
and the coefficients $\left\{  u_{lj}\right\}  $ form a $r\times r$ unitary
matrix i.e.
\begin{equation}
\mathbb{U}^{\dagger}\mathbb{U}=\mathbb{UU}^{\dagger}=\mathbb{I}_{r}.
\end{equation}


In Appendix \ref{Operators} we show that the idempotency condition
\begin{equation}
\left(  \sum_{1\leq l\leq r}\beta_{l}b_{l}^{\dagger}b_{l}\right)  ^{2}%
=\sum_{1\leq l\leq r}\beta_{l}b_{l}^{\dagger}b_{l}%
\end{equation}
together with the norm condition
\begin{equation}
\left\Vert \sum_{1\leq l\leq r}\beta_{l}b_{l}^{\dagger}b_{l}\right\Vert
_{\mathcal{B}\left(  \mathcal{H}_{\mathcal{F}}\right)  }=1
\end{equation}
can only be satisfied by positive operators where
\begin{equation}
\beta_{l}=0,\,1\leq l\neq k\leq r\text{ }%
\end{equation}
for some fixed $k.$ Thus all projectors belonging to $\mathcal{B}\left(
\mathcal{H}_{\mathcal{F}}\right)  _{1}$ are of the form
\begin{equation}
\lambda_{1}=b^{\dagger}b,
\end{equation}
where $b^{\dagger}$ creates a normalized General Spin Orbital (GSO)

\subsection{Projectors in $\mathcal{B}\left(  \mathcal{H}_{\mathcal{F}%
}\right)  _{2}$ \label{Proj2}}

Every positive operator in $\mathcal{B}\left(  \mathcal{H}_{\mathcal{F}%
}\right)  _{2}$ can be expressed as (Appendix \ref{Projection} )
\begin{equation}
\lambda_{2}=\sum_{\substack{1\leq j<k\leq r\\1\leq l<m\leq r}}\lambda
_{2jklm}a_{j}^{\dagger}a_{k}^{\dagger}a_{m}a_{l}=\sum_{1\leq l\leq\binom{r}%
{2}}\beta_{l}g_{l}^{\dagger}g_{l},\,\beta_{l}\geq0,1\leq l\leq\binom{r}{2},
\end{equation}
where the geminal creator is
\begin{equation}
g_{l}^{\dagger}=\sum_{1\leq j<k\leq r}u_{ljk}a_{j}^{\dagger}a_{k}^{\dagger
},\,1\leq l\leq\binom{r}{2}%
\end{equation}
and the coefficients $\left\{  \left.  u_{ljk}\right\vert 1\leq l\leq\binom
{r}{2},1\leq j<k\leq r\right\}  $ form a $\binom{r}{2}\times\binom{r}{2}$
unitary matrix i.e.
\begin{equation}
\mathbb{U}^{\dagger}\mathbb{U}=\mathbb{UU}^{\dagger}=\mathbb{I}_{\binom{r}{2}%
}.
\end{equation}
In Appendix \ref{Projection} we show that the idempotency condition
\begin{equation}
\left(  \sum_{1\leq l\leq\binom{r}{2}}\beta_{l}g_{l}^{\dagger}g_{l}\right)
^{2}=\sum_{1\leq l\leq\binom{r}{2}}\beta_{l}g_{l}^{\dagger}g_{l}%
\end{equation}
together with the norm condition
\begin{equation}
\left\Vert \sum_{1\leq l\leq\binom{r}{2}}\beta_{l}g_{l}^{\dagger}%
g_{l}\right\Vert _{\mathcal{B}\left(  \mathcal{H}_{\mathcal{F}}\right)  }=1
\end{equation}
can only be satisfied by positive operators where
\begin{equation}
\beta_{l}=0,\,1\leq l\neq k\leq\binom{r}{2}\text{ }%
\end{equation}
for some fixed $k$ and the rank of the $k^{th}$ geminal must be one. Thus in
order that a two electron operator $\lambda_{2}$ be a projection operator it
must be of the form
\begin{equation}
\lambda_{2}=a^{\dagger}b^{\dagger}ba,
\end{equation}
where $a^{\dagger}$ and $b^{\dagger}$ create normalized orthogonal GSO's.

\subsection{Projectors in $\Lambda_{2}$ \label{Proj1&2}}

\subsubsection{Projections from two hole operators}

By considering the positivity, idempotency and normalization relationship for
two \emph{hole} operators
\begin{equation}
\eta_{2}=\sum_{\substack{1\leq j<k\leq r\\1\leq l<m\leq r}}\eta_{2jklm}%
a_{j}a_{k}a_{m}^{\dagger}a_{l}^{\dagger}=\sum_{1\leq l\leq\binom{r}{2}}%
\gamma_{l}h_{l}h_{l}^{\dagger},\,\gamma_{l}\geq0,1\leq l\leq\binom{r}{2},
\end{equation}
where the geminal annihilator is
\begin{equation}
g_{l}=\sum_{1\leq j<k\leq r}u_{_{2}ljk}a_{j}a_{k},\,1\leq l\leq\binom{r}{2}%
\end{equation}
and the coefficients $\left\{  \left.  u_{2ljk}\right\vert 1\leq l\leq
\binom{r}{2},1\leq j<k\leq r\right\}  $ form a $\binom{r}{2}\times\binom{r}%
{2}$ unitary matrix i.e.
\begin{equation}
\mathbb{U}_{2}^{\dagger}\mathbb{U}_{2}=\mathbb{U}_{2}\mathbb{U}_{2}^{\dagger
}=\mathbb{I}_{\binom{r}{2}},
\end{equation}
one can carry out the same analysis as in subsection \ref{Proj2} and obtain
(Appendix \ref{Projection}) that in order that the two hole operator $\eta
_{2}$ be a projection operator it must be of the form
\begin{equation}
\eta_{2}=baa^{\dagger}b^{\dagger},
\end{equation}
where $a$ and $b$ annihilate normalized orthogonal GSO's. The projection
operator $\eta_{2}$ does not belong to $\mathcal{B}\left(  \mathcal{H}%
_{\mathcal{F}}\right)  _{2}$ but by using the canonical anticommutation
relationships eq. $\left(  \ref{CAR}\right)  $ one can observe that
\begin{equation}
\eta_{2}=baa^{\dagger}b^{\dagger}=I_{\mathcal{H}_{\mathcal{F}}}-a^{\dagger
}a-b^{\dagger}b+a^{\dagger}b^{\dagger}ba
\end{equation}
and thus
\begin{equation}
a^{\dagger}a+b^{\dagger}b-a^{\dagger}b^{\dagger}ba=I_{\mathcal{H}%
_{\mathcal{F}}}-\eta_{2},
\end{equation}
which shows that
\begin{equation}
\tilde{\lambda}_{2}=a^{\dagger}a+b^{\dagger}b-a^{\dagger}b^{\dagger}ba
\end{equation}
is a projection operator that belongs to $\Lambda_{2}.$

\subsubsection{Projections from products of one electron operators}

By considering the positivity, idempotency and normalization relationships for
products $\lambda_{1}\lambda_{1}^{\prime}$ of one electron operators of the
form eq. $\left(  \ref{1par}\right)  $ one can determine (Appendix
\ref{Projection}) that these products are projection operators when
\begin{equation}
\lambda_{1}\lambda_{1}^{\prime}=a^{\dagger}bb^{\dagger}a=a^{\dagger
}a-a^{\dagger}b^{\dagger}ba.
\end{equation}


One can prove (Appendix \ref{Projection}) that there are no other projection
operators of the form
\begin{equation}
\lambda=\sum_{1\leq j,k\leq r}\lambda_{1jk}a_{j}^{\dagger}a_{k}+\sum
_{\substack{1\leq j<k\leq r\\1\leq l<m\leq r}}\lambda_{2jklm}a_{j}^{\dagger
}a_{k}^{\dagger}a_{m}a_{l}%
\end{equation}
other than
\begin{equation}
a^{\dagger}a,a^{\dagger}b^{\dagger}ba,a^{\dagger}a-a^{\dagger}b^{\dagger
}ba,a^{\dagger}a+b^{\dagger}b-a^{\dagger}b^{\dagger}ba.
\end{equation}


\section{Generalized P, Q and G Conditions \label{PQG}}

The necessary P, Q and G conditions for representability of a $\binom{r}%
{2}\times\binom{r}{2}$ hermitian matrices were quickly discovered \cite{PQG}.
These conditions can be expressed in terms of the basic variables of the
ground state energy optimization problem formed by the complex matrix elements
$\left\{  \left.  D_{2j_{1}j_{2}k_{1}k_{2}}\right\vert 1\leq j_{1}<j_{2}\leq
r,1\leq k_{1}<k_{2}\leq r\right\}  $ of the matrix $\mathbb{D}_{2}$. In the
notation of this paper they are:

\subsection{P-Condition}

A potential Second Order Reduced Density Matrix (SORDM)\ matrix $\mathbb{D}%
_{2}$ given by
\begin{equation}
D_{2j_{1}j_{2}k_{1}k_{2}}=\left(  \left.  a_{k_{1}}^{\dagger}a_{k_{2}%
}^{\dagger}a_{j_{1}}a_{j_{2}}\right\vert D\right)  =\overline{D_{2k_{1}%
k_{2}j_{1}j_{2}}}%
\end{equation}
is semi definite i.e.
\begin{equation}
\mathbb{D}_{2}\geq0.
\end{equation}


\subsection{Q-Condition}

The matrix $\mathbb{Q}$ given by
\begin{equation}
Q_{j_{1}j_{2}k_{1}k_{2}}=\left(  \left.  a_{k_{1}}^{\dagger}a_{j_{1}}a_{k_{2}%
}^{\dagger}a_{j_{2}}\right\vert D\right)  =\overline{Q_{k_{1}k_{2}j_{1}j_{2}}}%
\end{equation}
is semi definite i.e.
\begin{equation}
\mathbb{Q}\geq0.
\end{equation}
If one invokes the pure $N$ condition (Section \ref{NRepCon}) this condition
becomes a necessary $N$-representability condition and the $\mathbb{Q}$ matrix
can be expressed in terms of the matrix $\mathbb{D}_{2}$ as
\begin{equation}
Q_{j_{1}j_{2}k_{1}k_{2}}=\frac{1}{\left(  N-1\right)  }\sum_{1\leq l\leq
r}D_{2j_{2}lk_{1}l}\delta_{j_{1}k_{2}}-D_{2j_{1}j_{2}k_{1}k_{2}}.
\end{equation}


\subsection{G-Condition}

The matrix $\mathbb{G}_{2}$ given by
\begin{equation}
G_{j_{1}j_{2}k_{1}k_{2}}=\left(  \left.  a_{j_{1}}a_{j_{2}}a_{k_{1}}^{\dagger
}a_{k_{2}}^{\dagger}\right\vert D\right)  =\overline{G_{2k_{1}k_{2}j_{1}j_{2}%
}}%
\end{equation}
is semi definite i.e.
\begin{equation}
\mathbb{G}\geq0.
\end{equation}
Again if one invokes the pure $N$ condition (Section \ref{NRepCon}) this
condition becomes a necessary $N$-representability condition and the
$\mathbb{G}$ matrix can be expressed in terms of the matrix $\mathbb{D}_{2}$
as
\begin{equation}
G_{j_{1}j_{2}k_{1}k_{2}}=\frac{1}{\left(  N-1\right)  }\sum_{1\leq l\leq
r}\left\{  D_{2j_{2}lk_{1}l}\delta_{j_{1}k_{2}}+D_{2j_{1}lk_{2}l}\delta
_{j_{2}k_{1}}\right\}  -D_{2j_{1}j_{2}k_{1}k_{2}}.
\end{equation}


The following generalized P, Q and G conditions, produced by examining the
values of the functional $f$ on projectors belonging to $\Lambda_{2},$ in
combination with the one electron conditions
\begin{equation}
1\geq f\left(  b^{\dagger}b\right)  \geq0\,\forall b,
\end{equation}
are \emph{Necessary and Sufficient }for the existence of positive extensions
of bounded linear functionals
\begin{equation}
f:\Lambda_{2}\rightarrow\mathbb{C}%
\end{equation}
to a state of $\mathcal{B}\left(  \mathcal{H}_{\mathcal{F}}\right)  $ and are
based on the theorems of Krien \cite{Krein}, Segal \cite{Segal} and Kadison
\cite{Kadison}.

\subsection{Generalized P-Condition}

A linear functional $f$ can be extended to a state if
\begin{equation}
1\geq f\left(  a^{\dagger}b^{\dagger}ba\right)  \geq0\,\forall a,b,
\end{equation}
where $\left\vert a\right\rangle $ and $\left\vert b\right\rangle $ are
normalized orthogonal GSO's belonging to $\mathcal{H}^{1}$ and $a^{\dagger}$
and $b^{\dagger}$ are the creators of those GSO's. As
\begin{equation}
D_{2abab}=f\left(  a^{\dagger}b^{\dagger}ba\right)
\end{equation}
this is a \emph{necessary }condition for ensemble representability.

\subsection{Generalized Q-Condition}

A linear functional $f$ can be extended to a state if
\begin{equation}
1\geq f\left(  a^{\dagger}a-a^{\dagger}b^{\dagger}ba\right)  \geq0\,\forall
a,b,
\end{equation}
which gives
\begin{equation}
1\geq D_{1aa}-D_{2abab}\geq0\,\forall a,b, \label{Q1}%
\end{equation}
where the diagonal elements of a possible First Order Reduced Density Matrix
(FORDM) are defined as
\begin{equation}
D_{1aa}=f\left(  a^{\dagger}a\right)  .
\end{equation}
Again this is a \emph{necessary }condition for ensemble representability.

\subsection{Generalized G-Condition}

A linear functional $f$ can be extended to a state if
\begin{equation}
1\geq f\left(  a^{\dagger}a+b^{\dagger}b-a^{\dagger}b^{\dagger}ba\right)
\geq0\,\forall a,b,
\end{equation}
which gives
\begin{equation}
1\geq D_{1aa}-D_{1bb}-D_{2abab}\geq0\,\forall a,b \label{G1}%
\end{equation}
and also a \emph{necessary }condition for ensemble representability.

\section{Electron and Pair Number Conditions \label{PartPair}}

In order to obtain $N$-representable extensions one must constrain the
functional $f$ to be dispersion free on the number operator i.e.
\begin{equation}
f\left(  \mathcal{N}_{1}^{2}\right)  =f\left(  \mathcal{N}_{1}\right)  ^{2},
\label{Dispfree}%
\end{equation}
where
\begin{equation}
\mathcal{N}_{1}=\sum_{1\leq j\leq r}a_{j}^{\dagger}a_{j}.
\end{equation}
One can show (Appendix \ref{Number}) that eq.$\left(  \ref{Dispfree}\right)  $
is equivalent to
\begin{equation}
f\left(  \mathcal{N}_{1}\right)  =N \label{NumPart}%
\end{equation}
and
\begin{equation}
f\left(  \mathcal{N}_{2}\right)  =\binom{N}{2} \label{Numpairs}%
\end{equation}
for some fixed $N,$ where the pair operator $\mathcal{N}_{2}$ is defined as
\begin{equation}
\mathcal{N}_{2}=\sum_{1\leq j<k\leq r}a_{j}^{\dagger}a_{k}^{\dagger}a_{k}%
a_{j}.
\end{equation}
If $f$ is dispersion free on the number operator the values of $f$ on
$\mathcal{B}\left(  \mathcal{H}_{\mathcal{F}}\right)  _{1}$ must be related to
its values on $\mathcal{B}\left(  \mathcal{H}_{\mathcal{F}}\right)  _{2}$ by
\begin{equation}
f\left(  a_{j}^{\dagger}a_{k}\right)  =\left(  N-1\right)  ^{-1}f\left(
a_{j}^{\dagger}\mathcal{N}_{1}a_{k}\right)  =\left(  N-1\right)  ^{-1}%
\sum_{1\leq l\leq r}f\left(  a_{j}^{\dagger}a_{l}^{\dagger}a_{l}a_{k}\right)
, \label{Def1D}%
\end{equation}
thus if the functional is only defined on $\mathcal{B}\left(  \mathcal{H}%
_{\mathcal{F}}\right)  _{2}$ one can define allowable extensions to
$\mathcal{B}\left(  \mathcal{H}_{\mathcal{F}}\right)  _{1}$ by eq. $\left(
\ref{Def1D}\right)  .$

In terms of matrix elements the definition eq. $\left(  \ref{Def1D}\right)  $
becomes
\begin{equation}
D_{1jk}=\left(  N-1\right)  ^{-1}\sum_{1\leq l\leq r}D_{2jlkl}%
\end{equation}
and one can see that if the value of the functional on the pair operator is
set by eq. $\left(  \ref{Numpairs}\right)  $ then the electron number
condition eq. $\left(  \ref{NumPart}\right)  $ is automatically ensured.

\section{N-Representability Conditions \label{NRepCon}}

The representability conditions in section \ref{PQG} can be complimented and
modified by the conditions in section \ref{PartPair} to form ensemble
\emph{Necessary and Sufficient }$N$-representability conditions which are:

\begin{enumerate}
\item Normalization%

\begin{equation}
\sum_{1\leq j<k\leq r}D_{2jkjk}=\binom{N}{2} \label{NormCon}%
\end{equation}


\item Generalized P-Condition%
\begin{equation}
1\geq D_{2abab}\geq0\,\forall a,b \label{PCond}%
\end{equation}


\item Generalized Q-Condition%
\begin{equation}
1\geq\left(  N-1\right)  ^{-1}\sum_{1\leq l\leq r}D_{2alal}-D_{2abab}%
\geq0\,\forall a,b \label{QCond}%
\end{equation}


\item Generalized G-Condition%
\begin{equation}
1\geq\left(  N-1\right)  ^{-1}\sum_{1\leq l\leq r}\left\{  D_{2alal}%
+D_{2blbl}\right\}  -D_{2abab}\geq0\,\forall a,b. \label{GCond}%
\end{equation}

\end{enumerate}

The conditions eqs. $\left(  \ref{NormCon},\ref{PCond},\ref{GCond}\right)  $
and the lower bound in eq. $\left(  \ref{QCond}\right)  $ have been discovered
in the study of the Diagonal Representability Problem
\cite{DavidsonI,DavidsonII,DavidsonIII,WeinholdWilson,Erdahl78}. However,
these conditions were not caste in terms of bounds on the extremal values of a
linear functional on the intersection of the unit ball of the Fermion Fock
algebra with the subspace $\Lambda_{2}$. Expressing the bounds in this context
allows one to utilize general extension theorems from functional analysis to
determine necessary and sufficient conditions for a linear functional to be a
reduced state \emph{if all the projectors in this intersection are
characterized}. The theorems in appendix \ref{Projection} characterize all
these projections.

\section{Discussion \label{Discussion}}

In this paper we have described how the $N$-Representability problem entails
finding the conditions that linear functionals belonging to the dual space,
$\mathcal{B}\left(  \mathcal{H}_{\mathcal{F}}\right)  _{2}^{\ast},$ of bounded
linear functionals defined on the subspace of second quantized two electron
operators, $\mathcal{B}\left(  \mathcal{H}_{\mathcal{F}}\right)  _{2},$
contained in the Fock algebra, $\mathcal{F},$ need to satisfy in order to be
extended to $N$-electron states and thus be reduced $N$-electron states i.e.
restrictions of $N$-electron states.

As this algebra is a $w^{\ast}$-algebra the theorems of Krein, Segal and
Kadison were utilized to show that the positive extension conditions could be
expressed in terms of the values of the restricted linear functionals on
projectors contained in the subspace $\Lambda_{2}$ generated by the one and
two electron operators, i.e. on the extreme points, $\operatorname*{ext}%
\left\{  \Lambda_{2+}\right\}  ,$ of the cone of positive operators belonging
to this space, which have to belong to the interval $\left[  0,1\right]  .$
These conditions are generalizations of the known P, Q and G conditions, which
are given by the lower bound of this interval. In addition the condition that
these functionals be restrictions of $N$-electron states entails that they
must be dispersion free on the number operator, $\mathcal{N}_{1},$ and that
their values on the one electron subspace $\mathcal{B}\left(  \mathcal{H}%
_{\mathcal{F}}\right)  _{1}$ be determined by their values on the two electron
subspace $\mathcal{B}\left(  \mathcal{H}_{\mathcal{F}}\right)  _{2}$.

We also described how the functionals belonging to $\mathcal{B}\left(
\mathcal{H}_{\mathcal{F}}\right)  _{2}^{\ast}$ correspond to operators
belonging $\mathcal{B}\left(  \mathcal{H}^{2}\right)  $ that act in
$\mathcal{H}^{2}$ and obtained their matrix representations in a natural way
using the generalized contraction map. This map is the adjoint of the
generalized expansion map which we display is the second quantization map.
This correspondence shows that the $N$-electron ground state of the system can
be determined by a constrained variation of the linear functional that
corresponds to the reduced Hamiltonian, over operators belonging to
$\mathcal{B}\left(  \mathcal{H}^{2}\right)  $ subject to the normalization and
generalized P, Q and G constraints, which are all linear. This problem can be
expressed in terms of matrix elements of operators belonging $\mathcal{B}%
\left(  \mathcal{H}^{2}\right)  $ produced by decomposable bases of
$\mathcal{H}^{2},$ however in order to be sufficient conditions the matrix
elements must satisfy them in \emph{all} \emph{decomposable} bases, i.e. the
conditions on the matrices produced by these elements must be conserved under
the under the action of the representations, $\mathcal{R}_{\mathcal{H}^{2}%
}\left(  U\left(  \mathcal{H}^{1}\right)  \right)  ,$ of the one electron
unitary group $U\left(  \mathcal{H}^{1}\right)  $ on $\mathcal{H}^{2}.$ These
observations can also be used to obtain to a more constrained form of the Semi
Definite Programming (SDP) method successfully used by Mazziotti
\cite{Mazziotti2011}.

In a future work this variational problem will be expressed in terms of
functions defined on the group manifold of $U\left(  \mathcal{H}^{1}\right)  $
and shown to lead to a classical Linear Programming problem that is amenable
to sequential approximations, that are all $N$-Representable.

\appendix


\section{Operator Spaces \label{Operators}}

In the following the Fock space $\mathcal{H}_{\mathcal{F}}$ will be considered
as the underlying Hilbert space on which the operators are defined and a few
basic properties will be listed.

\subsection{Finite Rank Operators}

An operator $X$ is a finite rank operator acting $\mathcal{H}_{\mathcal{F}}$
if
\begin{equation}
X=\sum_{1\leq j\leq m<\infty}\alpha_{j}\left\vert \Phi_{j}\right\rangle
\left\langle \Psi_{j}\right\vert ,
\end{equation}
where $\left\{  \left.  \Phi_{j}\right\vert 1\leq j\leq m\right\}  $ and
$\left\{  \left.  \Psi_{j}\right\vert 1\leq j\leq m\right\}  $ are orthonormal
sets in $\mathcal{H}_{\mathcal{F}}.$ The set of finite rank operators is
denoted $\mathcal{B}_{0}\left(  \mathcal{H}_{\mathcal{F}}\right)  $ and is not
closed if $\mathcal{H}_{\mathcal{F}}$ is infinite dimensional.

\subsection{Bounded Operators}

The operator norm (also called the uniform norm) is defined as
\begin{equation}
\left\Vert X\right\Vert _{\mathcal{B}\left(  \mathcal{H}_{\mathcal{F}}\right)
}=\sup_{\left\vert \Phi\right\rangle \in\mathcal{H}_{F}}\left\{  \left(
\frac{\left\langle \left.  \Phi\right\vert X^{\dagger}X\Phi\right\rangle
}{\left\langle \left.  \Phi\right\vert \Phi\right\rangle }\right)  ^{\frac
{1}{2}}\right\}  =\sup_{\left\vert \Phi\right\rangle \in\mathcal{H}_{F}%
}\left\{  \frac{\left\Vert X\Phi\right\Vert _{\mathcal{H}_{F}}}{\left\Vert
\Phi\right\Vert _{\mathcal{H}_{F}}}\right\}  ,
\end{equation}
(where $\left\Vert {}\right\Vert _{\mathcal{H}_{F}}$ is the norm in
$\mathcal{H}_{F}),$ and the space of bounded operators acting on the Hilbert
space $\mathcal{H}_{\mathcal{F}}$ as
\begin{equation}
\mathcal{B}\left(  \mathcal{H}_{\mathcal{F}}\right)  =\left\{  X\left\vert
\left\Vert X\right\Vert _{\mathcal{B}\left(  \mathcal{H}_{\mathcal{F}}\right)
}<\infty\right.  \right\}  .
\end{equation}
If one restricts attention to self adjoint operators $X=X^{\dagger}$ then it
can be shown that
\begin{equation}
\left\Vert X\right\Vert _{\mathcal{B}\left(  \mathcal{H}_{\mathcal{F}}\right)
}=\sup_{\left\vert \Phi\right\rangle \in\mathcal{H}_{F}}\left\{  \left\vert
\frac{\left\langle \left.  \Phi\right\vert X\Phi\right\rangle }{\left\langle
\left.  \Phi\right\vert \Phi\right\rangle }\right\vert \right\}
\end{equation}
and thus for positive operators, $X\geq0$%
\begin{equation}
\left\Vert X\right\Vert _{\mathcal{B}\left(  \mathcal{H}_{\mathcal{F}}\right)
}=\sup_{\left\vert \Phi\right\rangle \in\mathcal{H}_{F}}\left\{
\frac{\left\langle \left.  \Phi\right\vert X\Phi\right\rangle }{\left\langle
\left.  \Phi\right\vert \Phi\right\rangle }\right\}  .
\end{equation}
The space $\mathcal{B}\left(  \mathcal{H}_{\mathcal{F}}\right)  $ is closed in
the operator norm topology.

\subsection{Compact Operators}

The space of compact operators, $\mathcal{C}\left(  \mathcal{H}_{\mathcal{F}%
}\right)  ,$ is the closed space produced by taking the closure of the set of
finite rank operators with respect to the operator norm topology i.e. any
compact operator can be expressed as
\begin{equation}
X=\lim_{m\rightarrow\infty}\sum_{1\leq j\leq m}\alpha_{j}\left\vert \Phi
_{j}\right\rangle \left\langle \Psi_{j}\right\vert ,
\end{equation}
in the sense that
\begin{equation}
\lim_{m\rightarrow\infty}\left\Vert X-\sum_{1\leq j\leq m}\alpha_{j}\left\vert
\Phi_{j}\right\rangle \left\langle \Psi_{j}\right\vert \right\Vert
_{\mathcal{B}\left(  \mathcal{H}_{\mathcal{F}}\right)  }=0.
\end{equation}


\subsection{Trace Class Operators}

The trace class norm is defined as
\begin{equation}
\left\Vert X\right\Vert _{1}=\operatorname*{Tr}\left\{  \left(  X^{\dagger
}X\right)  ^{\frac{1}{2}}\right\}
\end{equation}
and the space of trace class operators acting on the Hilbert space
$\mathcal{H}_{\mathcal{F}}$ as
\begin{equation}
\mathcal{B}_{1}\left(  \mathcal{H}_{\mathcal{F}}\right)  =\left\{  X\left\vert
\left\Vert X\right\Vert _{1}<\infty\right.  \right\}  .
\end{equation}
The space $\mathcal{B}_{1}\left(  \mathcal{H}_{\mathcal{F}}\right)  $ is
closed with respect to the $\left\Vert {}\right\Vert _{1}$ norm, while the
closure with respect to the $\left\Vert {}\right\Vert _{\mathcal{B}\left(
\mathcal{H}_{\mathcal{F}}\right)  }$ norm produces $\mathcal{C}\left(
\mathcal{H}_{\mathcal{F}}\right)  .$

\subsection{Hilbert-Schmidt Operators}

The Hilbert-Schmidt norm is defined as
\begin{equation}
\left\Vert X\right\Vert _{2}=\left(  \operatorname*{Tr}\left\{  X^{\dagger
}X\right\}  \right)  ^{\frac{1}{2}}%
\end{equation}
and the space of Hilbert-Schmidt operators acting on the Hilbert space
$\mathcal{H}_{\mathcal{F}}$ as
\begin{equation}
\mathcal{B}_{2}\left(  \mathcal{H}_{\mathcal{F}}\right)  =\left\{  X\left\vert
\left\Vert X\right\Vert _{2}<\infty\right.  \right\}  .
\end{equation}
The space $\mathcal{B}_{2}\left(  \mathcal{H}_{\mathcal{F}}\right)  $ is
closed with respect to the $\left\Vert {}\right\Vert _{2}$ norm, while the
closure with respect to $\left\Vert {}\right\Vert _{\mathcal{B}\left(
\mathcal{H}_{\mathcal{F}}\right)  }$ norm produces $\mathcal{C}\left(
\mathcal{H}_{\mathcal{F}}\right)  .$

\subsection{Relationship between the Spaces}

The relationship between the spaces is
\begin{equation}
\mathcal{B}\left(  \mathcal{H}_{\mathcal{F}}\right)  \supseteq\mathcal{C}%
\left(  \mathcal{H}_{\mathcal{F}}\right)  \supseteq\mathcal{B}_{2}\left(
\mathcal{H}_{\mathcal{F}}\right)  \supseteq\mathcal{B}_{1}\left(
\mathcal{H}_{\mathcal{F}}\right)  \supseteq\mathcal{B}_{0}\left(
\mathcal{H}_{\mathcal{F}}\right)  .
\end{equation}
The closure of $\mathcal{B}_{0}\left(  \mathcal{H}_{\mathcal{F}}\right)  $
(and thus of $\mathcal{C}\left(  \mathcal{H}_{\mathcal{F}}\right)  ,$
$\mathcal{B}_{2}\left(  \mathcal{H}_{\mathcal{F}}\right)  $ and $\mathcal{B}%
_{1}\left(  \mathcal{H}_{\mathcal{F}}\right)  )$ with respect to the strong
operator topology is $\mathcal{B}\left(  \mathcal{H}_{\mathcal{F}}\right)  ,$
where these limits can be defined as
\begin{equation}
X=\text{\textrm{s-}}\lim_{m\rightarrow\infty}X_{m}\Leftrightarrow
\lim_{m\rightarrow\infty}\left\Vert \left(  X-X_{m}\right)  \Phi\right\Vert
_{\mathcal{H}_{\mathcal{F}}}=0\quad\forall\left\vert \Phi\right\rangle
\in\mathcal{H}_{\mathcal{F}},
\end{equation}
in contrast to the uniform operator topology limits that are given by%
\begin{align}
X  &  =\lim_{m\rightarrow\infty}X_{m}\Leftrightarrow\lim_{m\rightarrow\infty
}\left\Vert \left(  X-X_{m}\right)  \right\Vert _{\mathcal{H}_{\mathcal{F}}%
}=0\nonumber\\
&  =\lim_{m\rightarrow\infty}\sup_{\left\vert \Phi\right\rangle \in
\mathcal{H}_{F}}\left\{  \left(  \frac{\left\langle \left.  \Phi\right\vert
\left(  X-X_{m}\right)  ^{\dagger}\left(  X-X_{m}\right)  \Phi\right\rangle
}{\left\langle \left.  \Phi\right\vert \Phi\right\rangle }\right)  ^{\frac
{1}{2}}\right\}  .
\end{align}


If the dimension of $\mathcal{H}_{\mathcal{F}}$ is finite all of these spaces
are identical and closed in all of the preceding norms and topologies.

\section{Projectors \label{Projection}}

A general element of $\Lambda_{2}$ can be expanded as
\begin{equation}
\lambda=\sum_{1\leq j,k\leq r}\lambda_{1jk}a_{j}^{\dagger}a_{k}+\sum
_{\substack{1\leq j<k\leq r\\1\leq l<m\leq r}}\lambda_{2jklm}a_{j}^{\dagger
}a_{k}^{\dagger}a_{m}a_{l}%
\end{equation}
and the condition that it is a projector i.e. idempotent is that%
\begin{equation}
\lambda=\lambda^{2}, \label{idem}%
\end{equation}
which implies
\begin{equation}
\lambda=\lambda^{\dagger}\text{ and }\lambda\geq0.
\end{equation}
In a $c^{\ast}$-algebra
\begin{equation}
\left\Vert X^{\dagger}X\right\Vert =\left\Vert X\right\Vert ^{2} \label{Cstar}%
\end{equation}
The idempotency condition eq. $\left(  \ref{idem}\right)  $ implies that
\begin{equation}
\left\Vert \lambda^{2}\right\Vert =\left\Vert \lambda\right\Vert
\end{equation}
and the $c^{\ast}$-norm property eq. $\left(  \ref{Cstar}\right)  $%
\begin{equation}
\left\Vert \lambda^{2}\right\Vert =\left\Vert \lambda\right\Vert ^{2}%
\end{equation}
gives
\begin{equation}
\left\Vert \lambda\right\Vert ^{2}=\left\Vert \lambda\right\Vert
\end{equation}
and thus
\begin{equation}
\left\Vert \lambda\right\Vert =1.
\end{equation}


\subsection{One electron case}

We first consider the case $\lambda_{2}=0$ then
\begin{equation}
\lambda_{1}=\lambda_{1}^{2} \label{idem1}%
\end{equation}
and noting that $\lambda_{1}$ must be self adjoint and positive one has
\begin{equation}
\lambda=\lambda_{1}=\sum_{1\leq j,k\leq\infty}\lambda_{1jk}a_{j}^{\dagger
}a_{k}=\sum_{\substack{1\leq j,k\leq r\\1\leq l\leq r}}u_{jl}\beta_{l}\bar
{u}_{kl}a_{j}^{\dagger}a_{k}=\sum_{1\leq l\leq r}\beta_{l}b_{l}^{\dagger}%
b_{l},
\end{equation}
where%
\begin{subequations}
\begin{align}
1  &  \geq\beta_{l}\geq0;1\leq j\leq r\\
b_{l}^{\dagger}  &  =\sum_{1\leq j\leq r}a_{j}^{\dagger}u_{jl}\\
b_{l}^{\dagger}\left\vert \phi\right\rangle  &  =\left\vert \psi
_{l}\right\rangle \\
\mathbb{U}_{1}^{\dagger}\mathbb{U}_{1}  &  =\mathbb{U}_{1}\mathbb{U}%
_{1}^{\dagger}=\mathbb{I}_{r},
\end{align}
where $\mathbb{U}_{1}$ is the matrix composed of the matrix elements $\left\{
\left.  u_{jk}\right\vert 1\leq j,k\leq r\right\}  .$ The condition eq.
$\left(  \ref{idem1}\right)  $ becomes
\end{subequations}
\begin{equation}
\sum_{1\leq l\leq r}\beta_{l}b_{l}^{\dagger}b_{l}=\sum_{1\leq l\leq r}%
\beta_{l}^{2}b_{l}^{\dagger}b_{l}+2\sum_{1\leq l_{1}<l_{2}\leq r}\beta_{l_{1}%
}\beta_{l_{2}}b_{l_{1}}^{\dagger}b_{l_{1}}b_{l_{2}}^{\dagger}b_{l_{2}},
\end{equation}
which can be rewritten as
\begin{equation}
\sum_{1\leq l\leq r}\left\{  \beta_{l}\left(  1-\beta_{l}\right)  \right\}
b_{l}^{\dagger}b_{l}=2\sum_{1\leq l_{1}<l_{2}\leq r}\beta_{l_{1}}\beta_{l_{2}%
}b_{l_{1}}^{\dagger}b_{l_{2}}^{\dagger}b_{l_{2}}b_{l_{1}}.
\end{equation}
This equality must be true on all vectors in $\mathcal{H}_{\mathcal{F}}$ in
particular on $\left\vert \psi_{\kappa}\right\rangle \in\mathcal{H}^{1}%
,1\leq\kappa\leq r,$ i.e.%
\begin{equation}
\sum_{1\leq l\leq r}\left\{  \beta_{l}\left(  1-\beta_{l}\right)  \right\}
b_{l}^{\dagger}b_{l}\left\vert \psi_{\kappa}\right\rangle =2\sum_{1\leq
l_{1}<l_{2}\leq r}\beta_{l_{1}}\beta_{l_{2}}b_{l_{1}}^{\dagger}b_{l_{2}%
}^{\dagger}b_{l_{2}}b_{l_{1}}\left\vert \psi_{\kappa}\right\rangle ,
\end{equation}
which gives
\begin{equation}
\beta_{\kappa}\left(  1-\beta_{\kappa}\right)  \left\vert \psi_{\kappa
}\right\rangle =0,\quad1\leq\kappa\leq r
\end{equation}
i.e.
\begin{equation}
\beta_{\kappa}=0\text{ or }\beta_{\kappa}=1,\quad1\leq\kappa\leq r.\text{ }%
\end{equation}
One can see that $\beta_{\kappa}=1$ for only one value of $\kappa\in\left\{
1,\cdots,r\right\}  ,$ for let $\left\{  n_{1},\cdots,n_{r^{\prime}}\right\}
$ be a subset of $\left\{  1,\cdots,r\right\}  $ and
\begin{equation}
\lambda_{1}=\sum_{\substack{1\leq\nu\leq\kappa\leq r\\\sigma_{\nu}\in\left\{
1,\cdots,r\right\}  }}b_{\sigma_{\nu}}^{\dagger}b_{\sigma_{\nu}}%
\end{equation}
then
\begin{equation}
\lambda_{1}^{2}=\sum_{\sigma\in\left\{  1,\cdots,r\right\}  }b_{\sigma
}^{\dagger}b_{\sigma}+\sum_{\sigma,\sigma^{\prime}\in\left\{  1,\cdots
,r\right\}  }b_{\sigma}^{\dagger}b_{\sigma^{\prime}}^{\dagger}b_{\sigma
^{\prime}}b_{\sigma}=\lambda_{1}=\sum_{\sigma\in\left\{  1,\cdots,r\right\}
}b_{\sigma}^{\dagger}b_{\sigma}%
\end{equation}
thus
\begin{equation}
\sum_{\sigma,\sigma^{\prime}\in\left\{  1,\cdots,r\right\}  }b_{\sigma
}^{\dagger}b_{\sigma^{\prime}}^{\dagger}b_{\sigma^{\prime}}b_{\sigma}=0,
\end{equation}
which is only possible if $r^{\prime}=1.$ Thus if $\lambda_{1}$ is a projector
it must be of the form
\begin{equation}
\lambda_{1}=b^{\dagger}b.
\end{equation}


\subsection{Two electron case}

We consider the case $\lambda_{1}=0$ then
\begin{equation}
\lambda_{2}=\lambda_{2}^{2} \label{idem2}%
\end{equation}
and%
\begin{align}
\lambda &  =\lambda_{2}=\sum_{\substack{1\leq j_{1}<j_{2}\leq r\\1\leq
k_{1}<k_{2}\leq r}}\lambda_{2j_{1}j_{2}k_{1}k_{2}}a_{j_{1}}^{\dagger}a_{j_{2}%
}^{\dagger}a_{k_{2}}a_{k_{1}}=\sum_{\substack{1\leq j_{1}<j_{2}\leq r\\1\leq
k_{1}<k_{2}\leq r\\1\leq l\leq\binom{r}{2}}}u_{j_{1}j_{2}l}\beta_{l}\bar
{u}_{lk_{1}k_{2}}a_{j_{1}}^{\dagger}a_{j_{2}}^{\dagger}a_{k_{2}}a_{k_{1}%
}\nonumber\\
&  =\sum_{1\leq l\leq\binom{r}{2}}\beta_{l}g_{l}^{\dagger}g_{l},
\label{expan_lamda2}%
\end{align}
where the matrix that diagonalizes the matrix
\begin{equation}
\mathbf{\lambda}_{2}\equiv\left\{  \left.  \lambda_{2j_{1}j_{2}k_{1}k_{2}%
}\right\vert 1\leq j_{1}<j_{2}\leq r;1\leq k_{1}<k_{2}\leq r\right\}
\end{equation}
is unitary i.e.
\begin{equation}
\mathbb{U}_{2}^{\dagger}\mathbb{U}_{2}=\mathbb{U}_{2}\mathbb{U}_{2}^{\dagger
}=\mathbb{I}_{\binom{r}{2}},
\end{equation}
so that
\begin{equation}
\mathbf{\lambda}_{2}=\mathbb{U}_{2}\left(
\begin{array}
[c]{ccc}%
\beta_{1} & \cdots & 0\\
\vdots & \ddots & \vdots\\
0 & \cdots & \beta_{\binom{r}{2}}%
\end{array}
\right)  \mathbb{U}_{2}^{\dagger}.
\end{equation}
The geminal creators $\left\{  \left.  g_{l}^{\dagger}\right\vert 1\leq
l\leq\binom{r}{2}\right\}  $, geminal coefficients $\left\{  \left.
g_{lj_{1}j_{2}}\right\vert 1\leq j_{1}<j_{2}\leq r\right\}  $ and the geminal
coefficient matrices $\mathbb{G}_{l}$ have the properties%
\begin{subequations}
\begin{align}
&  g_{l}^{\dagger}=\sum_{1\leq j_{1}<j_{2}\leq r}g_{lj_{1}j_{2}}a_{j_{1}%
}^{\dagger}a_{j_{2}}^{\dagger}=\frac{1}{2}\sum_{1\leq j_{1},j_{2}\leq
r}g_{lj_{1}j_{2}}a_{j_{1}}^{\dagger}a_{j_{2}}^{\dagger};\,\,\\
&  \mathbb{G}_{lj_{1}j_{2}}=g_{lj_{1}j_{2}}=-g_{lj_{2}j_{1}}=u_{j_{1}j_{2}l}\\
&  g_{l}^{\dagger}\left\vert \phi\right\rangle =\left\vert g_{l}\right\rangle
\,;1\leq l\leq\binom{r}{2}\\
&  \left\langle g_{l_{1}}\left\vert g_{l_{2}}\right.  \right\rangle
=\delta_{l_{1}l_{2}};1\leq l_{1},l_{2}\leq\binom{r}{2}\\
&  \operatorname*{Tr}\left\{  \mathbb{G}_{l_{1}}^{\dagger}\mathbb{G}_{l_{2}%
}\right\}  =2\delta_{l_{1}l_{2}};1\leq l_{1},l_{2}\leq\binom{r}{2}\text{ }\\
&  \mathbb{G}_{l}=-\mathbb{G}_{l}^{t};1\leq l\leq\binom{r}{2}.
\end{align}
The values of the $\beta_{l}$'s in eq. $\left(  \ref{expan_lamda2}\right)  $
are restricted by
\end{subequations}
\begin{equation}
1\geq\beta_{l}\geq0;1\leq l\leq\binom{r}{2} \label{beta}%
\end{equation}
as the projection of $\lambda_{2}$ on $\mathcal{H}^{2}$ is given by
\begin{equation}
P_{\mathcal{H}^{2}}\lambda_{2}P_{\mathcal{H}^{2}}=\sum_{1\leq l\leq\binom
{r}{2}}\beta_{l}\left\vert g_{l}\right\rangle \left\langle g_{l}\right\vert ,
\end{equation}
where $\lambda_{2}$ is the second quantization of $P_{\mathcal{H}^{2}}%
\lambda_{2}P_{\mathcal{H}^{2}}$ i.e.%
\begin{equation}
\lambda_{2}=\Gamma\left(  P_{\mathcal{H}^{2}}\lambda_{2}P_{\mathcal{H}^{2}%
}\right)  ,
\end{equation}
thus%
\begin{equation}
\lambda_{2}\geq0\Rightarrow P_{\mathcal{H}^{2}}\lambda_{2}P_{\mathcal{H}^{2}%
}\geq0
\end{equation}
and%
\begin{equation}
\lambda_{2}-P_{\mathcal{H}^{2}}\lambda_{2}P_{\mathcal{H}^{2}}\geq0,
\end{equation}
as $\left\Vert \lambda_{2}\right\Vert =1$ ($\lambda_{2}$ is a projector) the
eigenvalues of $P_{\mathcal{H}^{2}}\lambda_{2}P_{\mathcal{H}^{2}}\in\left[
0,1\right]  .$

\subsubsection{Two Electron Operator Idempotency Condition}

\begin{lemma}
If $\lambda_{2}$ is a projector then $\lambda_{2}=\sum\limits_{1\leq
l\leq\binom{r}{2}}\beta_{l}g_{l}^{\dagger}g_{l},$ where $\beta_{l}=1$ or $0.$
\end{lemma}

\begin{proof}
Using the expansion eq. $\left(  \ref{expan_lamda2}\right)  $ of the operator
$\lambda_{2}$ the idempotency condition eq. $\left(  \ref{idem2}\right)  $
becomes
\begin{equation}
\sum_{1\leq l\leq\binom{r}{2}}\beta_{l}g_{l}^{\dagger}g_{l}=\sum_{1\leq
l_{1},l_{2}\leq\binom{r}{2}}\beta_{l_{1}}\beta_{l_{2}}g_{l_{1}}^{\dagger
}g_{l_{1}}g_{l_{2}}^{\dagger}g_{l_{2}} \label{lam2}%
\end{equation}
and one must compare the left hand side of eq. $\left(  \ref{lam2}\right)  $
to the right hand side. In order to do this one needs to bring the RHS of eq.
$\left(  \ref{lam2}\right)  $ to normal form by sequentially using the
canonical anti commutation relationships. The first step produces
\begin{align}
&  \sum_{1\leq l_{1},l_{2}\leq\binom{r}{2}}\beta_{l_{1}}\beta_{l_{2}}g_{l_{1}%
}^{\dagger}g_{l_{1}}g_{l_{2}}^{\dagger}g_{l_{2}}=\nonumber\\
&  \sum_{1\leq l_{1},l_{2}\leq\binom{r}{2}}\beta_{l_{1}}\beta_{l_{2}}\left\{
g_{l_{1}}^{\dagger}\left[  g_{l_{1}},g_{l_{2}}^{\dagger}\right]  g_{l_{2}%
}+g_{l_{1}}^{\dagger}g_{l_{2}}^{\dagger}g_{l_{1}}g_{l_{2}}\right\}  .
\label{normform}%
\end{align}
which can be developed by noting that%
\begin{equation}
\left[  g_{l_{1}},g_{l_{2}}^{\dagger}\right]  =\frac{1}{4}\sum
_{\substack{1\leq j_{1},j_{2}\leq r\\1\leq k_{1},k_{2}\leq r}}\bar{g}%
_{l_{1}j_{1}j_{2}}g_{l_{2}k_{1}k_{2}}\left[  a_{j_{2}}a_{j_{1}},a_{k_{1}%
}^{\dagger}a_{k_{2}}^{\dagger}\right]  =\,-\frac{1}{4}\sum_{\substack{1\leq
j_{1},j_{2}\leq r\\1\leq k_{1},k_{2}\leq r}}\bar{g}_{l_{1}j_{1}j_{2}}%
g_{l_{2}k_{1}k_{2}}\left[  a_{k_{1}}^{\dagger}a_{k_{2}}^{\dagger},a_{j_{2}%
}a_{j_{1}}\right]  ,
\end{equation}
and (Appendix \ref{Commutation Relationship})
\begin{align}
\left[  a_{k_{1}}^{\dagger}a_{k_{2}}^{\dagger},a_{j_{2}}a_{j_{1}}\right]   &
=I_{\mathcal{H}_{\mathcal{F}}}\delta_{j_{2}k_{1}}\delta_{j_{1}k_{2}%
}-I_{\mathcal{H}_{\mathcal{F}}}\delta_{j_{2}k_{2}}\delta_{j_{1}k_{1}}%
+a_{k_{2}}^{\dagger}a_{j_{2}}\delta_{j_{1}k_{1}}\nonumber\\
&  -a_{k_{1}}^{\dagger}a_{j_{2}}\delta_{j_{1}k_{2}}-a_{k_{2}}^{\dagger
}a_{j_{1}}\delta_{j_{2}k_{1}}+a_{k_{1}}^{\dagger}a_{j_{1}}\delta_{j_{2}k_{2}}%
\end{align}
to give
\begin{align}
&  \left[  g_{l_{1}},g_{l_{2}}^{\dagger}\right]  =-\frac{1}{4}\sum
_{\substack{1\leq j_{1},j_{2}\leq r\\1\leq k_{1},k_{2}\leq r}}\bar{g}%
_{l_{1}j_{1}j_{2}}g_{l_{2}k_{1}k_{2}}\left\{  I_{\mathcal{H}_{\mathcal{F}}%
}\delta_{j_{2}k_{1}}\delta_{j_{1}k_{2}}-I_{\mathcal{H}_{\mathcal{F}}}%
\delta_{j_{2}k_{2}}\delta_{j_{1}k_{1}}\right. \nonumber\\
&  \left.  \qquad\qquad\qquad\qquad\qquad\qquad\right.  \left.  a_{k_{2}%
}^{\dagger}a_{j_{2}}\delta_{j_{1}k_{1}}-a_{k_{1}}^{\dagger}a_{j_{2}}%
\delta_{j_{1}k_{2}}-a_{k_{2}}^{\dagger}a_{j_{1}}\delta_{j_{2}k_{1}}+a_{k_{1}%
}^{\dagger}a_{j_{1}}\delta_{j_{2}k_{2}}\right\}  ,\nonumber\\
&  =\frac{1}{2}\sum_{1\leq j_{1},j_{2}\leq r}\bar{g}_{l_{1}j_{1}j_{2}}%
g_{l_{2}j_{1}j_{2}}\nonumber\\
&  \qquad\qquad\qquad-\frac{1}{4}\sum_{\substack{1\leq j,j_{2}\leq r\\1\leq
k_{2}\leq r}}\left\{  \bar{g}_{l_{1}jj_{2}}g_{l_{2}jk_{2}}-\bar{g}%
_{l_{1}jj_{2}}g_{l_{2}k_{2}j}-\bar{g}_{l_{1}j_{2}j}g_{l_{2}jk_{2}}+\bar
{g}_{l_{1}j_{2}j}g_{l_{2}k_{2}j}\right\}  a_{k_{2}}^{\dagger}a_{j_{2}},
\end{align}
and after using $g_{ljk}=-g_{lkj}$
\begin{equation}
\left[  g_{l_{1}},g_{l_{2}}^{\dagger}\right]  =\delta_{l_{1}l_{2}}+\sum_{1\leq
j,k\leq r}\bar{g}_{l_{1}jj}g_{l_{2}jk}a_{k}^{\dagger}a_{j}=\delta_{l_{1}l_{2}%
}-\sum_{1\leq j_{2},k_{2}\leq r}\left(  \mathbb{G}_{l_{1}}^{\dagger}%
\mathbb{G}_{l_{2}}\right)  _{jk}a_{k}^{\dagger}a_{j}.
\end{equation}
The idempotency condition then becomes%
\begin{align}
&  \sum_{1\leq l\leq\binom{r}{2}}\beta_{l}g_{l}^{\dagger}g_{l}=\sum_{1\leq
l_{1},l_{2}\leq\binom{r}{2}}\beta_{l_{1}}\beta_{l_{2}}g_{l_{1}}^{\dagger
}g_{l_{1}}g_{l_{2}}^{\dagger}g_{l_{2}}\nonumber\\
&  =\sum_{1\leq l_{1},l_{2}\leq\binom{r}{2}}\beta_{l_{1}}\beta_{l_{2}}\left\{
g_{l_{1}}^{\dagger}\left\{  \delta_{l_{1}l_{2}}-\sum_{1\leq j_{2},k_{2}\leq
r}\left(  \mathbb{G}_{l_{1}}^{\dagger}\mathbb{G}_{l_{2}}\right)  _{j_{2}k_{2}%
}a_{k_{2}}^{\dagger}a_{j_{2}}\right\}  g_{l_{2}}+g_{l_{1}}^{\dagger}g_{l_{2}%
}^{\dagger}g_{l_{1}}g_{l_{2}}\right\}  ,
\end{align}
which simplifies to%
\begin{align}
&  \sum_{1\leq l\leq\binom{r}{2}}\beta_{l}g_{l}^{\dagger}g_{l}=\nonumber\\
&  =\sum_{1\leq l\leq\binom{r}{2}}\beta_{l}^{2}g_{l}^{\dagger}g_{l}%
-\sum_{1\leq l_{1},l_{2}\leq\binom{r}{2}}\beta_{l_{1}}\beta_{l_{2}}\left\{
\sum_{1\leq j_{2},k_{2}\leq r}\left(  \mathbb{G}_{l_{1}}^{\dagger}%
\mathbb{G}_{l_{2}}\right)  _{j_{2}k_{2}}g_{l_{1}}^{\dagger}a_{k_{2}}^{\dagger
}a_{j_{2}}g_{l_{2}}+g_{l_{1}}^{\dagger}g_{l_{2}}^{\dagger}g_{l_{1}}g_{l_{2}%
}\right\} \nonumber\\
&  =\sum_{1\leq l\leq\binom{r}{2}}\beta_{l}^{2}g_{l}^{\dagger}g_{l}%
-\sum_{1\leq l_{1},l_{2}\leq\binom{r}{2}}\beta_{l_{1}}\beta_{l_{2}}\sum_{1\leq
j_{2},k_{2}\leq r}\left(  \mathbb{G}_{l_{1}}^{\dagger}\mathbb{G}_{l_{2}%
}\right)  _{j_{2}k_{2}}g_{l_{1}}^{\dagger}a_{k_{2}}^{\dagger}a_{j_{2}}%
g_{l_{2}}\nonumber\\
&  \qquad\qquad\qquad\qquad\qquad\qquad\qquad\qquad\qquad\qquad\qquad
\qquad-\sum_{1\leq l_{1},l_{2}\leq\binom{r}{2}}\beta_{l_{1}}\beta_{l_{2}%
}g_{l_{1}}^{\dagger}g_{l_{2}}^{\dagger}g_{l_{1}}g_{l_{2}}. \label{work}%
\end{align}
The operator equality in eq. $\left(  \ref{work}\right)  $ must be true on all
vectors thus in particular
\begin{equation}
\sum_{1\leq l\leq\binom{r}{2}}\beta_{l}g_{l}^{\dagger}g_{l}\left\vert
g_{p}\right\rangle =\sum_{1\leq l\leq\binom{r}{2}}\beta_{l}^{2}g_{l}^{\dagger
}g_{l}\left\vert g_{p}\right\rangle \Longleftrightarrow\beta_{p}\left(
\beta_{p}-1\right)  =0;\quad1\leq p\leq\binom{r}{2}%
\end{equation}
leading to
\begin{equation}
\lambda_{2}=\sum_{\substack{1\leq\nu\leq\kappa\leq\binom{r}{2}\\\sigma_{\nu
}\in\left\{  1,\cdots,\binom{r}{2}\right\}  }}g_{\sigma_{\nu}}^{\dagger
}g_{\sigma_{\nu}};\quad\sigma_{\nu}\neq\sigma_{\nu^{\prime}}\text{ unless }%
\nu=\nu^{\prime},
\end{equation}
where $\kappa$ is the number of non zero $\beta^{\prime}s.$
\end{proof}

\begin{lemma}
If $\lambda_{2}$ is a projector then $\lambda_{2}=g^{\dagger}g,$ where
$g^{\dagger}$ creates a normalized geminal.
\end{lemma}

\begin{proof}
As $\lambda_{2}$ is a projection operator its eigenvalues are $1$ or $0$ and
as $\left[  \lambda_{2},\mathcal{N}_{1}\right]  =0$ its eigenvectors can be
chosen to simultaneously diagonalize the number operator $\mathcal{N}_{1}.$

In the finite dimensional case, i.e. $r<\infty,$ the subspace $\mathcal{H}%
^{r}$ is one dimensional and is spanned by a single normalized r-fold\emph{
}antisymmetrized product, $\left\vert \varphi_{1}\cdots\varphi_{r}%
\right\rangle ,$ formed from \emph{any }orthonormal basis, $\left\{  \left.
\varphi_{j}\right\vert 1\leq j\leq r\right\}  $ of $\mathcal{H}^{1},$ i.e.
\begin{equation}
\left\vert \varphi_{1}\cdots\varphi_{r}\right\rangle \overset{\text{up to a
phase factor}}{=}\left\vert \varphi_{1}^{\prime}\cdots\varphi_{r}^{\prime
}\right\rangle ,
\end{equation}
where $\left\{  \left.  \varphi_{j}\right\vert 1\leq j\leq r\right\}  $ and
$\left\{  \left.  \varphi_{j}^{\prime}\right\vert 1\leq j\leq r\right\}  $ are
arbitrary orthonormal bases of $\mathcal{H}^{1}.$Thus any vector in
$\mathcal{H}^{r}$ is an eigenvector of $\lambda_{2},$ as $\lambda_{2}:$
$\mathcal{H}^{r}\rightarrow\mathcal{H}^{r}$.

Any geminal creator $g_{l}^{\dagger}$ can be expressed in canonical form as
\cite{OrtizWeiner2006,Mazziotti2001}
\begin{equation}
g_{l}^{\dagger}=\sum_{1\leq j\leq s}c_{lj}a_{lj}^{\dagger}a_{lj+s}^{\dagger
};\,c_{lj}\geq0,1\leq j\leq s,
\end{equation}
where
\begin{equation}
\sum_{1\leq j\leq s}c_{lj}^{2}=1
\end{equation}
and the canonical general spin orbitals of $\left\vert g_{l}\right\rangle =$
$g_{l}^{\dagger}\left\vert \phi\right\rangle $ are produced by
\begin{equation}
a_{lj}^{\dagger}\left\vert \phi\right\rangle =\left\vert \varphi
_{lj}\right\rangle .
\end{equation}
Hence
\begin{equation}
g_{l}^{\dagger}g_{l}=\sum_{1\leq j,k\leq s}c_{lj}c_{lk}a_{lj}^{\dagger
}a_{lj+s}^{\dagger}a_{lk+s}a_{lk}%
\end{equation}
and $\lambda_{2}$ can be expanded (noting that $\beta_{l}=1$ or $0,$ $1\leq
l\leq\binom{r}{2}$ ) as
\begin{align}
&  \left\langle \varphi_{1}\cdots\varphi_{r}\left\vert \lambda_{2}\varphi
_{1}\cdots\varphi_{r}\right.  \right\rangle =\sum_{1\leq l\leq\binom{r}{2}%
}\beta_{l}\left\langle \varphi_{l1}\cdots\varphi_{lr}\left\vert g_{l}%
^{\dagger}g_{l}\varphi_{l1}\cdots\varphi_{lr}\right.  \right\rangle \\
&  =\sum_{\substack{1\leq j,k\leq s\\1\leq l\leq\binom{r}{2}}}\beta_{l}%
c_{lj}c_{lk}\left\langle \varphi_{l1}\cdots\varphi_{lr}\left\vert
a_{lj}^{\dagger}a_{lj+s}^{\dagger}a_{lk+s}a_{lk}\varphi_{l1}\cdots\varphi
_{lr}\right.  \right\rangle \\
&  =\sum_{\substack{1\leq j\leq s\\1\leq l\leq\binom{r}{2}}}\beta_{l}%
c_{lj}^{2}=\sum_{1\leq l\leq\binom{r}{2}}\beta_{l}.
\end{align}
The vector $\left\vert \varphi_{1}\cdots\varphi_{r}\right\rangle $ is an
eigenvector of $\lambda_{2}$ with eigenvalue $1$ (as $\lambda_{2}$ is a
projector) thus
\begin{equation}
\sum_{1\leq l\leq\binom{r}{2}}\beta_{l}=1,
\end{equation}
as $\beta_{l}=1$ or $0,$ $1\leq l\leq\binom{r}{2}$ only one $\beta_{l}$ can be
non zero and $\lambda=g^{\dagger}g$.
\end{proof}

Thus one can prove the following theorem:

\begin{theorem}
\bigskip$\lambda_{2}\in\mathcal{B}\left(  \mathcal{H}_{\mathcal{F}}\right)
_{2}$ is a projector if only if it can be expressed in the form $\lambda
_{2}=a^{\dagger}b^{\dagger}ba,$ where $a^{\dagger}$ and $b^{\dagger}$ create
normalized orthogonal vectors in $\mathcal{H}^{1}.$
\end{theorem}

\begin{proof}
Using the two preceding lemmas one has that if $\lambda_{2}$ is a two electron
projector
\begin{equation}
\lambda_{2}=g^{\dagger}g\nonumber
\end{equation}
and thus
\begin{equation}
\lambda_{2}^{2}=g^{\dagger}gg^{\dagger}g. \label{ProjSq}%
\end{equation}
Using the canonical expansion of the geminal one can expand eq. $\left(
\ref{ProjSq}\right)  $ as
\begin{equation}
g^{\dagger}gg^{\dagger}g=\sum_{1\leq j_{1},j_{2},j_{3},j_{4}\leq s}a_{j_{1}%
}^{\dagger}a_{j_{1+s}}^{\dagger}a_{j_{2}}a_{j_{2+s}}a_{j_{3}}^{\dagger
}a_{j_{3+s}}^{\dagger}a_{j_{4}}a_{j_{4+s}}c_{j_{1}}c_{j_{2}}c_{j_{3}}c_{j_{4}%
}. \label{Projsq2}%
\end{equation}
and bring eq. $\left(  \ref{Projsq2}\right)  $ to normal form by observing
that
\begin{align}
&  a_{j_{2}}a_{j_{2+s}}a_{j_{3}}^{\dagger}a_{j_{3+s}}^{\dagger}\nonumber\\
&  =\delta_{j_{2}j_{3+s}}\delta_{j_{2+s}j_{3}}-\delta_{j_{2}j_{3}}%
\delta_{j_{2+s}j_{3+s}}-a_{j_{3+s}}^{\dagger}a_{j_{2}}\delta_{j_{2+s}j_{3}%
}+\nonumber\\
&  \quad+a_{j_{3}}^{\dagger}a_{j_{2}}\delta_{j_{2+s}j_{3+s}}+a_{j_{3+s}%
}^{\dagger}a_{j_{2+s}}\delta_{j_{2}j_{3}}-a_{j_{3}}^{\dagger}a_{j_{2+s}}%
\delta_{j_{2}j_{3+s}}+a_{j_{3}}^{\dagger}a_{j_{3+s}}^{\dagger}a_{j_{2+s}%
}a_{j_{2}}\nonumber\\
&  =-I\delta_{j_{2}j_{3}}+a_{j_{3}}^{\dagger}a_{j_{2}}\delta_{j_{2}j_{3}%
}+a_{j_{3}}^{\dagger}a_{j_{2}}\delta_{j_{2}j_{3}}+a_{j_{3}}^{\dagger
}a_{j_{3+s}}^{\dagger}a_{j_{2+s}}a_{j_{2}}%
\end{align}
and obtaining%
\begin{align}
&  g^{\dagger}gg^{\dagger}g=\sum_{1\leq j_{1},j_{2},j_{3},j_{4}\leq s}%
a_{j_{1}}^{\dagger}a_{j_{1+s}}^{\dagger}\left(  -\delta_{j_{2}j_{3}}+a_{j_{3}%
}^{\dagger}a_{j_{2}}\delta_{j_{2}j_{3}}+a_{j_{3}+s}^{\dagger}a_{j_{2}+s}%
\delta_{j_{2}j_{3}}+a_{j_{3}}^{\dagger}a_{j_{3+s}}^{\dagger}a_{j_{2+s}%
}a_{j_{2}}\right) \nonumber\\
&  \qquad\qquad\qquad\times a_{j_{4}}a_{j_{4+s}}c_{j_{1}}c_{j_{2}}c_{j_{3}%
}c_{j_{4}}\nonumber\\
&  =\sum_{1\leq j_{1},j_{4}\leq s}a_{j_{1}}^{\dagger}a_{j_{1+s}}^{\dagger
}a_{j_{4+s}}a_{j_{4}}c_{j_{1}}c_{j_{4}}\left\Vert g\right\Vert _{\mathcal{H}%
^{2}}^{2}\nonumber\\
&  \qquad+\sum_{1\leq j_{1},j_{2},j_{4}\leq s}c_{j_{1}}a_{j_{1}}^{\dagger
}a_{j_{1+s}}^{\dagger}n_{j_{2}}\left(  a_{j_{2}}^{\dagger}a_{j_{2}}%
+a_{j_{2}+s}^{\dagger}a_{j_{2}+s}\right)  a_{j_{4}}a_{j_{4+s}}c_{j_{4}%
}\nonumber\\
&  \quad\qquad+\sum_{1\leq j_{1},j_{2},j_{3},j_{4}\leq s}a_{j_{1}}^{\dagger
}a_{j_{1+s}}^{\dagger}a_{j_{3}}^{\dagger}a_{j_{3+s}}^{\dagger}a_{j_{2+s}%
}a_{j_{2}}a_{j_{4}}a_{j_{4+s}}c_{j_{1}}c_{j_{2}}c_{j_{3}}c_{j_{4}},
\end{align}
where $n_{j}=a_{j}^{\dagger}a_{j}.$ The difference between $\lambda_{2}^{2}$
and $\lambda_{2}$ can then be expressed as
\begin{align}
\Xi &  =g^{\dagger}gg^{\dagger}g-g^{\dagger}g=\sum_{1\leq j_{1},j_{2}%
,j_{4}\leq s}c_{j_{1}}a_{j_{1}}^{\dagger}a_{j_{1+s}}^{\dagger}n_{j_{2}}\left(
a_{j_{2}}^{\dagger}a_{j_{2}}+a_{j_{2}+s}^{\dagger}a_{j_{2}+s}\right)
a_{j_{4}}a_{j_{4+s}}c_{j_{4}}\nonumber\\
&  +\sum_{1\leq j_{1},j_{2},j_{3},j_{4}\leq s}a_{j_{1}}^{\dagger}a_{j_{1+s}%
}^{\dagger}a_{j_{3}}^{\dagger}a_{j_{3+s}}^{\dagger}a_{j_{2+s}}a_{j_{2}%
}a_{j_{4}}a_{j_{4+s}}c_{j_{1}}c_{j_{2}}c_{j_{3}}c_{j_{4}}%
\end{align}
and one can note that
\begin{gather}
\Xi=g^{\dagger}gg^{\dagger}g-g^{\dagger}g=0\nonumber\\
\Updownarrow\nonumber\\
P_{\mathcal{H}^{m}}\Xi P_{\mathcal{H}^{m}}=0\forall m.
\end{gather}
Thus in particular
\begin{align}
\left\langle \Phi\left\vert \sum_{1\leq j_{1},j_{2},j_{4}\leq s}c_{j_{1}%
}a_{j_{1}}^{\dagger}a_{j_{1+s}}^{\dagger}n_{j_{2}}\left(  a_{j_{2}}^{\dagger
}a_{j_{2}}+a_{j_{2}+s}^{\dagger}a_{j_{2}+s}\right)  a_{j_{4}}a_{j_{4+s}%
}c_{j_{4}}\right.  \Psi\right\rangle  &  =0\nonumber\\
\forall\Phi,\Psi &  \in\mathcal{H}^{3}.
\end{align}
Letting
\begin{equation}
\left\vert \Phi\right\rangle =\left\vert \Psi\right\rangle =\left\vert
\varphi_{j}\varphi_{j+s}\varphi_{k}\right\rangle
\end{equation}
one obtains
\begin{align}
0  &  =\left\langle \varphi_{j}\varphi_{j+s}\varphi_{k}\left\vert \sum_{1\leq
j_{1},j_{2},j_{4}\leq s}c_{j_{1}}a_{j_{1}}^{\dagger}a_{j_{1+s}}^{\dagger
}n_{j_{2}}\left(  a_{j_{2}}^{\dagger}a_{j_{2}}+a_{j_{2}+s}^{\dagger}%
a_{j_{2}+s}\right)  a_{j_{4}}a_{j_{4+s}}c_{j_{4}}\varphi_{j}\varphi
_{j+s}\varphi_{k}\right.  \right\rangle \nonumber\\
&  =\sum_{1\leq j_{1},j_{2},j_{4}\leq s}c_{j_{1}}c_{j_{4}}n_{j_{2}}\left\{
\left\langle \varphi_{j}\varphi_{j+s}\varphi_{k}\left\vert a_{j_{1}}^{\dagger
}a_{j_{1+s}}^{\dagger}a_{j_{2}}^{\dagger}a_{j_{2}}a_{j_{4}}a_{j_{4+s}}%
\varphi_{j}\varphi_{j+s}\varphi_{k}\right.  \right\rangle \right. \nonumber\\
&  \left.  +\right.  \left.  \left\langle \varphi_{j}\varphi_{j+s}\varphi
_{k}\left\vert a_{j_{1}}^{\dagger}a_{j_{1+s}}^{\dagger}a_{j_{2}+s}^{\dagger
}a_{j_{2}+s}a_{j_{4}}a_{j_{4+s}}\varphi_{j}\varphi_{j+s}\varphi_{k}\right.
\right\rangle \right\}  =n_{j}n_{k}\left(  1-\delta_{jk}-\delta_{\left(
j+s\right)  k}\right)  .
\end{align}
Thus
\begin{equation}
n_{j}n_{\left\vert k\right\vert }=0\quad j\neq k\text{ };1\leq j\leq s,1\leq
k\leq2s,
\end{equation}
i.e.
\begin{align}
n_{1}n_{2}  &  =n_{1}n_{3}=\cdots=n_{1}n_{s}=0\nonumber\\
n_{2}n_{1}  &  =n_{2}n_{3}=\cdots=n_{2}n_{s}=0, \label{Proof}%
\end{align}
where
\begin{equation}
\left\vert k\right\vert =\left\{
\begin{array}
[c]{c}%
k\text{ if }k\leq s\\
k-s\text{ if }k>s
\end{array}
\right.  .
\end{equation}
One can then see from eq. $\left(  \ref{Proof}\right)  $ that all but one of
the occupation numbers must be zero and the one that isn't must equal one,
thus the rank of the geminal creator $g^{\dagger}$ must be $1.$
\end{proof}

\subsection{Mixed Case}

The general case concerns projectors of the form
\begin{equation}
\lambda=\lambda_{1}+\lambda_{2},
\end{equation}
which is handled by considering:

\subsubsection{Q Projectors}

Using the particle-hole transformation
\begin{equation}
a^{\dagger}\rightarrow a
\end{equation}
of the Fock algebra and the preceding lemmas and theorem one can see that all
projectors formed from two electron hole operators must be of the form%
\begin{equation}
A^{\dagger}A,
\end{equation}
where
\begin{equation}
A=b^{\dagger}a^{\dagger},
\end{equation}
i.e.%
\begin{equation}
A^{\dagger}A=abb^{\dagger}a^{\dagger},
\end{equation}
which is a projector as
\begin{align}
&  \left(  abb^{\dagger}a^{\dagger}\right)  ^{2}=abb^{\dagger}a^{\dagger
}abb^{\dagger}a^{\dagger}=abb^{\dagger}\left(  1-aa^{\dagger}\right)
bb^{\dagger}a^{\dagger}=\nonumber\\
&  =abb^{\dagger}bb^{\dagger}a^{\dagger}-abb^{\dagger}aa^{\dagger}bb^{\dagger
}a^{\dagger}=ab\left(  1-bb^{\dagger}\right)  b^{\dagger}a^{\dagger
}=abb^{\dagger}a^{\dagger}.
\end{align}
Noting that
\begin{equation}
abb^{\dagger}a^{\dagger}=I_{\mathcal{H}_{\mathcal{F}}}-a^{\dagger}%
a-b^{\dagger}b+a^{\dagger}b^{\dagger}ba=I_{\mathcal{H}_{\mathcal{F}}}-\Lambda,
\end{equation}
where
\begin{equation}
\Lambda=a^{\dagger}a+b^{\dagger}b-a^{\dagger}b^{\dagger}ba
\end{equation}
one has
\begin{equation}
I_{\mathcal{H}_{\mathcal{F}}}-\Lambda=\left(  I_{\mathcal{H}_{\mathcal{F}}%
}-\Lambda\right)  ^{2}=I_{\mathcal{H}_{\mathcal{F}}}-2\Lambda+\Lambda^{2},
\end{equation}
which shows that
\begin{equation}
\Lambda^{2}=\Lambda,
\end{equation}
i.e.
\begin{equation}
\Lambda=a^{\dagger}a+b^{\dagger}b-a^{\dagger}b^{\dagger}ba
\end{equation}
are projectors that belong to $\mathcal{B}\left(  \mathcal{H}_{\mathcal{F}%
}\right)  _{1}\oplus\mathcal{B}\left(  \mathcal{H}_{\mathcal{F}}\right)  _{2}$
and lead to the generalized Q-conditions.

\subsubsection{G Projectors}

The generalized G-condition is based on the product
\begin{equation}
A^{\dagger}A,
\end{equation}
where
\begin{equation}
A=b^{\dagger}a,
\end{equation}
i.e.%
\begin{equation}
A^{\dagger}A=a^{\dagger}bb^{\dagger}a=a^{\dagger}\left(  I_{\mathcal{H}%
_{\mathcal{F}}}-b^{\dagger}a\right)  a=a^{\dagger}a-a^{\dagger}b^{\dagger
}ba=n_{a}-n_{a}n_{b}=n_{a}\left(  1-n_{b}\right)  .
\end{equation}
The operator $a^{\dagger}a-a^{\dagger}b^{\dagger}ba$ is a projector as
\begin{equation}
n_{a}\left(  1-n_{b}\right)  n_{a}\left(  1-n_{b}\right)  =n_{a}\left(
1-n_{b}\right)  ^{2}=n_{a}\left(  1-2n_{b}+n_{b}\right)  =n_{a}\left(
1-n_{b}\right)
\end{equation}
and
\begin{equation}
a^{\dagger}a-a^{\dagger}b^{\dagger}ba\in\Lambda_{2}.
\end{equation}


\section{Two Electron Commutation Relationships
\label{Commutation Relationship}}

The commutator%
\begin{equation}
\left[  a_{k_{1}}^{\dagger}a_{k_{2}}^{\dagger},a_{j_{2}}a_{j_{1}}\right]
\end{equation}
is evaluated by bringing $a_{j_{4}}a_{j_{3}}a_{j_{1}}^{\dagger}a_{j_{2}%
}^{\dagger}$ to normal form as%
\begin{align}
&  a_{j_{2}}a_{j_{1}}a_{k_{1}}^{\dagger}a_{k_{2}}^{\dagger}=a_{j_{2}}\left(
\delta_{j_{1}k_{1}}-a_{k_{1}}^{\dagger}a_{j_{1}}\right)  a_{k_{2}}^{\dagger
}=a_{j_{2}}a_{k_{2}}^{\dagger}\delta_{j_{1}k_{1}}-a_{j_{2}}a_{k_{1}}^{\dagger
}a_{j_{1}}a_{k_{2}}^{\dagger}\nonumber\\
&  =\left(  \delta_{j_{2}k_{2}}-a_{k_{2}}^{\dagger}a_{j_{2}}\right)
\delta_{j_{1}k_{1}}-\left(  \delta_{j_{2}k_{1}}-a_{k_{1}}^{\dagger}a_{j_{2}%
}\right)  \left(  \delta_{j_{1}k_{2}}-a_{k_{2}}^{\dagger}a_{j_{1}}\right)
\nonumber\\
&  =\delta_{j_{2}k_{2}}\delta_{j_{1}k_{1}}-a_{k_{2}}^{\dagger}a_{j_{2}}%
\delta_{j_{1}k_{1}}-\delta_{j_{2}k_{1}}\delta_{j_{1}k_{2}}+a_{k_{1}}^{\dagger
}a_{j_{2}}\delta_{j_{1}k_{2}}+a_{k_{2}}^{\dagger}a_{j_{1}}\delta_{j_{2}k_{1}%
}-a_{k_{1}}^{\dagger}a_{j_{2}}a_{k_{2}}^{\dagger}a_{j_{1}}\nonumber\\
&  =\delta_{j_{2}k_{2}}\delta_{j_{1}k_{1}}-\delta_{j_{2}k_{1}}\delta
_{j_{1}k_{2}}-a_{k_{2}}^{\dagger}a_{j_{2}}\delta_{j_{1}k_{1}}+a_{k_{1}%
}^{\dagger}a_{j_{2}}\delta_{j_{1}k_{2}}+a_{k_{2}}^{\dagger}a_{j_{1}}%
\delta_{j_{2}k_{1}}-a_{k_{1}}^{\dagger}\left(  \delta_{j_{2}k_{2}}-a_{k_{2}%
}^{\dagger}a_{j_{2}}\right)  a_{j_{1}}\nonumber\\
&  =\delta_{j_{2}k_{2}}\delta_{j_{1}k_{1}}-\delta_{j_{2}k_{1}}\delta
_{j_{1}k_{2}}-a_{k_{2}}^{\dagger}a_{j_{2}}\delta_{j_{1}k_{1}}+a_{k_{1}%
}^{\dagger}a_{j_{2}}\delta_{j_{1}k_{2}}+a_{k_{2}}^{\dagger}a_{j_{1}}%
\delta_{j_{2}k_{1}}-a_{k_{1}}^{\dagger}a_{j_{1}}\delta_{j_{2}k_{2}}+a_{k_{1}%
}^{\dagger}a_{k_{2}}^{\dagger}a_{j_{2}}a_{j_{1}}%
\end{align}
giving
\begin{align}
&  a_{j_{2}}a_{j_{1}}a_{k_{1}}^{\dagger}a_{k_{2}}^{\dagger}=\nonumber\\
I_{\mathcal{H}_{\mathcal{F}}}  &  \delta_{j_{2}k_{2}}\delta_{j_{1}k_{1}%
}-I_{\mathcal{H}_{\mathcal{F}}}\delta_{j_{2}k_{1}}\delta_{j_{1}k_{2}}%
-a_{k_{2}}^{\dagger}a_{j_{2}}\delta_{j_{1}k_{1}}+a_{k_{1}}^{\dagger}a_{j_{2}%
}\delta_{j_{1}k_{2}}+a_{k_{2}}^{\dagger}a_{j_{1}}\delta_{j_{2}k_{1}}-a_{k_{1}%
}^{\dagger}a_{j_{1}}\delta_{j_{2}k_{2}}+a_{k_{1}}^{\dagger}a_{k_{2}}^{\dagger
}a_{j_{2}}a_{j_{1}},
\end{align}
and thus
\begin{equation}
\left[  a_{k_{1}}^{\dagger}a_{k_{2}}^{\dagger},a_{j_{2}}a_{j_{1}}\right]
=I_{\mathcal{H}_{\mathcal{F}}}\delta_{j_{2}k_{1}}\delta_{j_{1}k_{2}%
}-I_{\mathcal{H}_{\mathcal{F}}}\delta_{j_{2}k_{2}}\delta_{j_{1}k_{1}}%
+a_{k_{2}}^{\dagger}a_{j_{2}}\delta_{j_{1}k_{1}}-a_{k_{1}}^{\dagger}a_{j_{2}%
}\delta_{j_{1}k_{2}}-a_{k_{2}}^{\dagger}a_{j_{1}}\delta_{j_{2}k_{1}}+a_{k_{1}%
}^{\dagger}a_{j_{1}}\delta_{j_{2}k_{2}}.
\end{equation}


\section{Second Quantization Map \label{SecQuant}}

By considering the restrictions of $a_{j}^{\dagger}$ to $\mathcal{H}^{N}$,
i.e. $a_{j}^{\dagger}P_{\mathcal{H}^{N}},$ for $0\leq N\leq\infty$ one obtains%
\begin{align}
a_{j}^{\dagger}  &  =\sum_{0\leq N\leq\infty}\sum_{1\leq j_{1}<\cdots
<j_{N}\leq\infty}\left\vert \varphi_{j}\varphi_{j_{1}}\cdots\varphi_{j_{N}%
}\right\rangle \left\langle \varphi_{j_{1}}\cdots\varphi_{j_{N}}\right\vert
\nonumber\\
&  =\sum_{0\leq N\leq\infty}\sum_{1\leq j_{1},\cdots,j_{N}\leq\infty}%
\sqrt{\left(  N+1\right)  }\left\vert \varphi_{j}\wedge\varphi_{j_{1}}%
\wedge\cdots\wedge\varphi_{j_{N}}\right\rangle \left\langle \varphi_{j_{1}%
}\wedge\cdots\wedge\varphi_{j_{N}}\right\vert .
\end{align}
thus one can define the second quantization map $\Gamma$
\begin{equation}
\Gamma:\mathcal{F\rightarrow F}%
\end{equation}
from the Fock algebra to itself by
\begin{align}
\Gamma\left(  \left\vert \varphi_{j}\right\rangle \left\langle \phi\right\vert
\right)   &  =a_{j}^{\dagger}=\sum_{0\leq N\leq\infty}\sum_{1\leq j_{1}%
<\cdots<j_{N}\leq\infty}\left\vert \varphi_{j}\varphi_{j_{1}}\cdots
\varphi_{j_{N}}\right\rangle \left\langle \varphi_{j_{1}}\cdots\varphi_{j_{N}%
}\right\vert \nonumber\\
&  =\sum_{0\leq N\leq\infty}\sum_{1\leq j_{1},\cdots,j_{N}\leq\infty}%
\sqrt{\left(  N+1\right)  }\left\vert \varphi_{j}\wedge\varphi_{j_{1}}%
\wedge\cdots\wedge\varphi_{j_{N}}\right\rangle \left\langle \varphi_{j_{1}%
}\wedge\cdots\wedge\varphi_{j_{N}}\right\vert \nonumber\\
&  =\sum_{1\leq j_{1},\cdots,j_{N}\leq\infty}\sqrt{\left(  N+1\right)
}\left\vert \varphi_{j}\right\rangle \left\langle \phi\right\vert \wedge
P_{\mathcal{H}^{N}}.
\end{align}
Annihilators can be expressed in an analogous fashion by
\begin{align}
a_{j}  &  =\left\vert \phi\right\rangle \left\langle \varphi_{j}\right\vert
+\sum_{1\leq N\leq\infty}\sum_{1\leq j_{1}<\cdots<j_{N}\leq\infty}\left\vert
\varphi_{j_{1}}\cdots\varphi_{j_{N}}\right\rangle \left\langle \varphi
_{j}\varphi_{j_{1}}\cdots\varphi_{j_{N}}\right\vert \nonumber\\
&  =\sum_{0\leq N\leq\infty}\sum_{1\leq j_{1}<\cdots<j_{N}\leq\infty
}\left\vert \varphi_{j_{1}}\cdots\varphi_{j_{N}}\right\rangle \left\langle
\varphi_{j}\varphi_{j_{1}}\cdots\varphi_{j_{N}}\right\vert .
\end{align}
and%
\begin{equation}
a_{j}=\Gamma\left(  \left\vert \phi\right\rangle \left\langle \varphi
_{j}\right\vert \right)  =\sum_{1\leq j_{1},\cdots,j_{N}\leq\infty}%
\sqrt{\left(  N+1\right)  }\left\vert \phi\right\rangle \left\langle
\varphi_{j}\right\vert \wedge P_{\mathcal{H}^{N}}.\nonumber
\end{equation}
In a similar fashion
\begin{align}
&  a_{l_{1}}^{\dagger}a_{m_{1}}=\Gamma\left(  \left\vert \varphi_{l_{1}%
}\right\rangle \left\langle \varphi_{m_{1}}\right\vert \right) \nonumber\\
&  =\sum_{\substack{0\leq N,M\leq\infty\\1\leq j_{1}<\cdots<j_{N}\leq
\infty\\1\leq k_{1}<\cdots<k_{M}\leq\infty}}\left\vert \varphi_{l_{1}}%
\varphi_{j_{1}}\cdots\varphi_{j_{N}}\right\rangle \left\langle \left.
\varphi_{j_{1}}\cdots\varphi_{j_{N}}\right\vert \varphi_{k_{1}}\cdots
\varphi_{k_{M}}\right\rangle \left\langle \varphi_{m_{1}}\varphi_{k_{1}}%
\cdots\varphi_{k_{M}}\right\vert \nonumber\\
&  =\sum_{1\leq N\leq\infty}\sum_{1\leq j_{1}<\cdots<j_{N}\leq\infty
}\left\vert \varphi_{l_{1}}\varphi_{j_{1}}\cdots\varphi_{j_{N}}\right\rangle
\left\langle \varphi_{m_{1}}\varphi_{j_{1}}\cdots\varphi_{j_{N}}\right\vert
=\sum_{0\leq N\leq\infty}\left(  N+1\right)  \left\vert \varphi_{l_{1}%
}\right\rangle \left\langle \varphi_{m_{1}}\right\vert \wedge P_{\mathcal{H}%
^{N}}.
\end{align}
and%
\begin{align}
&  a_{l_{1}}^{\dagger}a_{l_{2}}^{\dagger}a_{m_{2}}a_{m_{1}}=\Gamma\left(
\left\vert \varphi_{l_{1}}\varphi_{l_{2}}\right\rangle \left\langle
\varphi_{m_{1}}\varphi_{m_{2}}\right\vert \right) \nonumber\\
&  =\sum_{\substack{0\leq N,M\leq\infty\\1\leq j_{1}<\cdots<j_{N}\leq
\infty\\1\leq k_{1}<\cdots<k_{M}\leq\infty}}\left\vert \varphi_{l_{1}}%
\varphi_{l_{2}}\varphi_{j_{1}}\cdots\varphi_{j_{N}}\right\rangle \left\langle
\left.  \varphi_{j_{1}}\cdots\varphi_{j_{N}}\right\vert \varphi_{k_{1}}%
\cdots\varphi_{k_{M}}\right\rangle \left\langle \varphi_{m_{1}}\varphi_{m_{2}%
}\varphi_{k_{1}}\cdots\varphi_{k_{M}}\right\vert \nonumber\\
&  =\sum_{1\leq N\leq\infty}\sum_{1\leq j_{1}<\cdots<j_{N}\leq\infty
}\left\vert \varphi_{l_{1}}\varphi_{l_{2}}\varphi_{j_{1}}\cdots\varphi_{j_{N}%
}\right\rangle \left\langle \varphi_{m_{1}}\varphi_{m_{2}}\varphi_{j_{1}%
}\cdots\varphi_{j_{N}}\right\vert =\sum_{0\leq N\leq\infty}\binom{N}%
{2}\left\vert \varphi_{l_{1}}\varphi_{l_{2}}\right\rangle \left\langle
\varphi_{m_{1}}\varphi_{m_{2}}\right\vert \wedge P_{\mathcal{H}^{N}}.
\end{align}
One \emph{should note} that in general
\begin{equation}
\Gamma\left(  A\right)  \Gamma\left(  B\right)  \neq\Gamma\left(  A\wedge
B\right)  \neq\Gamma\left(  AB\right)  ,
\end{equation}
so $\Gamma$ is not an algebraic map.

\section{States of a Fixed Number of Electrons \label{Number}}

\begin{theorem}
A state $f$ is a $N$ electron state iff $f\left(  \mathcal{N}_{1}\right)  =N$
and $f\left(  \mathcal{N}_{2}\right)  =\binom{N}{2}$ i.e. $f\left(
\mathcal{N}_{1}\right)  =N$ and $f\left(  \mathcal{N}_{2}\right)  =\binom
{N}{2}\Leftrightarrow f\left(  \left(  \mathcal{N}_{1}\right)  ^{2}\right)
=f\left(  \mathcal{N}_{1}\right)  ^{2}=N^{2}$
\end{theorem}

\begin{proof}
In the following we use the operator identity $\left(  \mathcal{N}_{1}\right)
^{2}=\mathcal{N}_{1}+2\mathcal{N}_{2}.$

If
\begin{equation}
f\left(  \mathcal{N}_{1}\right)  =N\text{ and }f\left(  \mathcal{N}%
_{2}\right)  =\binom{N}{2}%
\end{equation}
then%
\begin{equation}
f\left(  \mathcal{N}_{1}+2\mathcal{N}_{2}\right)  =f\left(  \mathcal{N}%
_{1}^{2}\right)  =N+\frac{2N\left(  N-1\right)  }{2}=N^{2}=f\left(
\mathcal{N}_{1}\right)  ^{2}%
\end{equation}
thus $f$ is a pure $N$ -electron state. The only if part is trivial as
$f\left(  \mathcal{N}_{1}\right)  ^{2}=N^{2}$ clearly implies that $f\left(
\mathcal{N}_{1}\right)  =N$ and $f\left(  \mathcal{N}_{2}\right)  =\binom
{N}{2}.$
\end{proof}

\bibliographystyle{apsrev}
\bibliography{PQGconditions}



\end{document}